% Thesis draft — the angle (anticipatory call-uncertainty around undisclosed cash deals)
% Structure: supervisor template (docs/Thesis/DraftTemplate.txt) — one-column (user order 2026-06-09).
% RULES: docs/Draft/thesis_tables.tex is the ONE reference document for every table and number.
%        Tables are embedded via _tables_from_bible.tex, generated BYTE-EXACT from the bible by
%        tmp/extract_draft_tables.py — never hand-edited. Regenerate after any bible change.
%        Numbers were extracted/verified programmatically (tmp/verify_draft_numbers.py — ALL CHECKS PASS, 2026-06-09).
%        Prose units are user-ratified one by one before they enter this file.
\documentclass[11pt]{article}
\usepackage[margin=2.5cm]{geometry}
\usepackage{newtxtext,newtxmath}
\usepackage{booktabs}
\usepackage{amsmath}
\usepackage{graphicx}
\usepackage[authoryear]{natbib}
\usepackage[hidelinks]{hyperref}
\pagestyle{plain}

\title{Cash Got Your Tongue? CEO Language Uncertainty Around Undisclosed Cash Acquisitions}
\author{Sina Soleimanipour\thanks{I am deeply grateful to my co-advisors, Dr.~Ali Akyol and Dr.~Harshit Rajaiya, whose guidance, rigorous feedback, and steady support shaped every part of this thesis. All remaining errors are my own.}\\[2pt]\small Telfer School of Management, University of Ottawa \\ \small ssole077@uottawa.ca}
\date{June 2026}

\begin{document}

\maketitle
\begin{center}
\begin{minipage}{0.85\textwidth}
\textbf{Abstract.} We document an anticipatory, cash-deal-concentrated component in CEO earnings-call language. Using the Dzielinski--Wagner--Zeckhauser decomposition to isolate the call-varying residual of Q\&A uncertainty for US public firms (2002--2018), we find that this residual uncertainty rises in the quarter immediately before an undisclosed cash acquisition (0.0461, $p<.01$), with no comparable positive run-up before stock-financed deals, and subsides at announcement even as the acquirer's accumulated cash persists to completion --- uncertainty follows the information clock, cash the transaction clock. The cash--stock difference survives a formal test, and analyst scrutiny of cash does not account for the run-up. The pattern is correlational; we read it as call language tracking a deal's disclosure state, and leave the underlying mechanism open.\par\vspace{4pt}
\noindent\textbf{Keywords:} earnings calls; CEO language; textual analysis; mergers and acquisitions; cash holdings; disclosure timing\par
\noindent\textbf{JEL Classification:} G14, G32, G34, D82
\end{minipage}
\vspace{18pt}
\end{center}

\section{Introduction}\label{sec:intro}

Between the decision to acquire another firm and the public announcement of that decision, a chief executive must keep doing something ordinary: hosting the quarterly earnings call. While the deal is pending it is material nonpublic information, so securities law forbids discussing it; the call, however, is a standing obligation that cannot be skipped, and much of it --- the question-and-answer session --- is unscripted, leaving the CEO to field analyst questions that may press on exactly what cannot be said. We show that this bind leaves a measurable trace. In the quarter immediately before a firm announces an undisclosed cash acquisition, the uncertainty in its CEO's unscripted answers --- the portion specific to that call rather than to the executive's habitual way of speaking --- runs higher than in the same firm's ordinary quarters, and by the quarter the deal becomes public it is gone. The cash the acquirer has been accumulating to pay for the purchase behaves differently, persisting until the deal closes.

The stakes are practical. If anticipation leaves a linguistic trace, then a routine earnings call carries information about a material event the firm has not yet disclosed --- information legible before any announcement makes the event public. A growing literature already studies corporate disclosure around acquisitions, though along dimensions other than this one: bidders manage the tone of their earnings press releases ahead of stock-for-stock deals \citep{thewissen2024}, the volume of corporate-strategy vocabulary on conference calls rises with deal activity \citep{ragozzino2024}, recent work examines the precision of management's guidance around transactions \citep{everhart2025}, and markets react to acquisition plans once those plans are disclosed \citep{gokkaya2025}. What this work has not measured, to our knowledge, is uncertainty --- whether the unscripted answers of a CEO privately committed to a deal grow measurably less certain before the market knows anything is afoot. We extend this lens rather than contest it, and the extension runs along three axes: from the tone and volume of language to its residual uncertainty, from stock-financed deals to cash-financed ones, and from the effects that follow an announcement to the anticipatory window that precedes it.

Our research question is deliberately descriptive: how does the uncertainty in a CEO's earnings-call language behave around undisclosed acquisitions --- when it rises, where it concentrates, and when it subsides? We ask how the pattern looks, not what produces it; throughout, we claim no identification and read the evidence as correlational. Now is the moment to ask the question. The two papers closest to ours both appeared in 2024, so the lens that reads deal information out of corporate language is being assembled as we write; its uncertainty dimension is still missing, and the measurement needed to capture it has only recently become tractable.

% dwz bib entry (names verified by user 2026-06-09): Dzielinski, Wagner, and Zeckhauser,
% "Straight talkers and vague talkers: The effects of managerial style in earnings conference calls."
We take the question to the earnings calls of US public firms from 2002 to 2018, linking the Capital IQ transcript archive --- parsed so that the CEO's unscripted answers are kept separate from the scripted presentation --- to acquisition records from SDC that date each deal's announcement and completion and record whether it was paid in cash or in stock. The measure of unscripted uncertainty follows the decomposition of \citet{dwz}: we regress the CEO's question-and-answer uncertainty on the executive's own fixed effect and on the call's circumstances and keep the residual, the call-specific innovation that neither the CEO's persistent style nor the firm's situation explains. That residual is the only place an anticipatory, deal-tracking signal could live, since such a signal must appear and disappear within a single executive's tenure rather than persist as a trait. Three designs then put the prediction to work: a within-firm comparison of a cash acquirer's pre-announcement quarter with its own ordinary quarters, an event study that follows the residual against the firm's cash across the stages of the deal's life --- before announcement, announced but not yet closed, and completed --- and a pooled test of whether the pre-announcement rise is larger for cash deals than for stock.

Three findings anchor the paper. First, residual question-and-answer uncertainty rises in the quarter immediately preceding the announcement of a cash acquisition (0.0461, significant at the 1\% level; Table~\ref{tab:empire_building_did}); among stock-financed acquirers, who labor under the same disclosure constraint but carry no comparable cash build-up, there is no such positive run-up ($-0.0429$, n.s.; same table). Second, the elevation follows the deal's clock rather than the balance sheet's: in a matched event window it is statistically indistinguishable from zero as soon as the deal is announced, even as the acquirer's accumulated cash stays on the books until the transaction closes (Table~\ref{tab:empire_drop_matched}). Third, the gap between the cash and stock arms survives a formal, pooled test of their difference (0.0983, significant at the 5\% level; Table~\ref{tab:empire_cashspec}), though, with the stock-side estimate imprecise, we treat it as supportive rather than definitive. A natural alternative --- that pre-deal quarters simply draw harder analyst questions about cash, and that the hard questions, not the deal, produce the hedging --- does not account for the run-up (Table~\ref{tab:reason_gating}). What the paper documents is the timing, the concentration in cash, and the resolution at announcement; the reason the pattern arises --- a CEO bound by disclosure law or one choosing strategic silence --- it leaves as a deliberately hedged interpretation.

The rest of the paper proceeds as follows. Section~\ref{sec:framework} develops the disclosure-state framework, states the predictions it implies, defines the residual measure, and sets out the empirical design. Section~\ref{sec:main} describes the data and sample and presents the three main analyses --- the run-up, its timing around the announcement, and its concentration in cash. Section~\ref{sec:additional} reports the rule-out of the analyst-scrutiny alternative and the contrast with the scripted presentation, and Section~\ref{sec:conclusion} concludes.

\section{Conceptual Framework and Empirical Strategy}\label{sec:framework}

\subsection{Conceptual Framework}


A firm that has privately committed to an acquisition but not yet announced it occupies a distinct \emph{disclosure state} rather than a distinct balance-sheet state. A pending acquisition that is material to the acquirer constitutes material nonpublic information: information about preliminary merger negotiations can be material well before any definitive agreement, with materiality determined by balancing the probability that the transaction is consummated against its magnitude to the firm \citep{basic1988}. Securities law imposes no general duty to disclose such confidential negotiations, so silence is permitted; but once a firm speaks it may not mislead, since an untrue or half-true statement is unlawful \citep{rule10b5, basic1988}. The quarterly earnings call is a recurring event that managers nonetheless hold, so the executive faces live questioning while unable to address the one material development and unable to deny it. Disclosure theory tells us such a withholding state is not informationally empty: an informed manager may rationally withhold private information when disclosure is costly, sustaining a threshold below which information is withheld \citep{verrecchia1983}, and non-disclosure persists in equilibrium precisely because outside investors cannot distinguish a manager who is uninformed from one who is informed but silent \citep{dye1985}. This withholding bind---material information held back while the firm keeps fielding questions---is the organizing primitive of our analysis, which we read in descriptive, correlational terms throughout.

The venue in which this bind becomes observable is the quarterly earnings conference call, which pairs a scripted managerial presentation with an unscripted question-and-answer session in which analysts question management directly. The two segments are not equally revealing: both are incrementally informative relative to the accompanying press release, and the discussion carries information content beyond the presentation \citep{matsumoto2011}---making the unscripted exchange, which managers cannot fully prepare in advance, the natural place to observe a constrained executive under questioning. That language can be read quantitatively: because general-purpose word lists misclassify financial text---nearly three-quarters of the Harvard dictionary's negative-word counts are attributable to words that are typically not negative in a financial context---\citet{lm2011} construct dictionaries specific to financial disclosure, among them an uncertainty list of 285 words organized around imprecision rather than risk. We adopt this established apparatus---the presentation/Q\&A split and the finance-specific dictionary, including its uncertainty list---to read the language of the call.

Within this venue the bind has an observable consequence. Because the pending acquisition is material and nonpublic, the executive can neither address the analysts' questions that bear on it nor deny that anything is afoot, so a direct answer is unavailable. Such non-answers are common, and they are not informationally empty: in roughly six of ten calls managers withhold requested information, leaving questions unanswered, and, in the words of that literature, silence speaks \citep{hollander2010}. But silence is not a costless refuge---precisely because it speaks---and the setting is an unscripted exchange rather than prepared remarks. Non-disclosure therefore need not surface as outright silence; it can surface in the words themselves, as answers that grow hedged, qualified, and non-committal---elevated uncertainty language in the Q\&A. We read this language, rather than the incidence of silence, because the constraint binds while the executive keeps speaking. The resulting pattern is keyed to the disclosure state: the bind predicts that uncertainty language is elevated while the acquisition remains undisclosed and recedes once announcement makes it public---a directional prediction we test rather than assume. This anticipatory timing---present before the announcement, absent after---is the first of the two dimensions we develop.

The anticipatory pattern just described, however, is not the only thing moving uncertainty language on a call. How uncertain a CEO sounds is partly a durable personal trait: some executives simply speak more cautiously than others, quarter after quarter. \citet{bertrand_schoar2003} show that managers carry such persistent, individual styles---manager fixed effects---that durably shape the policies of the firms they run, with a manager's imprint in one firm strongly predicting his imprint in the next across investment, financial, and organizational policy alike. \citet{dwz} construct such a decomposition for the language of earnings calls, separating a CEO's use of uncertainty words into a persistent, manager-specific style component and a time-varying, call-level residual. The implication for our purpose is a locating one. A persistent style is a fixed trait of the speaker, not a signal that tracks a particular, still-undisclosed transaction; an anticipatory component, by contrast, must emerge while a given deal is pending and subside once it is announced---it must vary from call to call within a single executive's tenure. If such a component exists, it therefore cannot reside in the persistent style; it can only live in the call-varying residual that remains once persistent style is netted out. This locates where an anticipatory signal must be---it does not claim the residual is that signal, which only the research design that follows can isolate.

A second dimension concerns not the timing of the signal but the deals in which it should concentrate. Cash and stock acquisitions share the identical disclosure setting---both are material, both are withheld---yet differ in one design-relevant respect: a cash purchase draws on an accumulated cash position, whereas a stock exchange need not. \citet{harford1999} documents the underlying fact: firms accumulate cash reserves above a baseline model of normal holdings, and cash-rich firms are more likely to make acquisitions. That accumulation is a by-product of retained free cash flow rather than saving toward a particular deal, so we take from it only that a cash bid draws on a pre-existing cash position---not that firms stockpile in order to fund planned acquisitions. This makes the stock deal a natural placebo: the same withholding bind, minus the cash commitment. A signal that concentrates in cash deals rather than stock deals is therefore informative on its own terms, even though why the two should differ is a question we leave open. The two dimensions also run on different clocks: the anticipatory signal tracks the undisclosed information and resolves once the deal is announced, while the cash commitment serves the purchase itself and---mechanically---persists until it is paid at completion. The contrast is one of timing: the uncertainty is resolved by the announcement, the cash is not.

This paper's nearest neighbors read deal-related language along dimensions other than ours. \citet{thewissen2024} show that acquirers manage the tone of their earnings press releases: stock bidders inflate disclosure tone by about 15 percent in the year before a stock-for-stock announcement, lifting their share price to reduce the cost of the shares they must issue. \citet{ragozzino2024} find a parallel pattern in the spoken record---on the earnings calls of firms active in M\&A, the volume of corporate-strategy vocabulary rises, with analysts using such terms nearly 9 percent more, and executives 7.2 percent more, than at inactive firms. Both study language that is deliberately deployed or topic-driven, and both concern tone or subject matter rather than expressed uncertainty. Separately, \citet{keown1981} document that prices often move ahead of a deal: roughly half of the market's reaction to a merger occurs before its first public announcement. That anticipatory signal, however, lives in prices, not in language. The cell these literatures leave empty is the one this paper occupies: uncertainty language in the unscripted Q\&A, in the anticipatory window before the deal is public, read for what it may reveal about a deal the firm is withholding rather than as language crafted to move the market. To our knowledge, no prior work reads this; we offer it as a positioning claim about where the contribution sits, not as a tested mechanism.

Two readings of this prediction are observationally equivalent, and the framework does not require us to choose between them. Under the first---compliance-constrained disclosure---the executive is unable to confirm the pending deal and barred from denying it, is left to answer around it, and uncertainty rises as a by-product of the bind. Under the second---strategic silence---the reticence is deliberate: the executive chooses to stay vague about a development that is not yet public. Both would produce the same observable signature: elevated uncertainty language in the Q\&A while the deal is undisclosed, resolving at announcement. Our design cannot distinguish them, and we do not try. What we characterize is therefore a correlational pattern keyed to the disclosure state---an anticipatory, cash-concentrated component of CEO uncertainty language that tracks a deal's passage from private to public---and not a mechanism. We identify no causal channel, establish no intent, and read the relationship descriptively throughout. It is in this deliberately bounded register that the empirical strategy of the following sections puts the pattern to the test.

\subsection{Hypothesis Development}

The framework yields three predictions.

\textbf{H1 (run-up).} While a deal is pending but undisclosed, the gag binds: residual Q\&A uncertainty should be elevated for a cash acquirer in the quarter immediately before announcement, relative to its own typical quarters.

\textbf{H1a (cash concentration).} The bind should be stronger where the balance sheet invites the forbidden question --- cash-financed deals --- than in stock-financed deals, which sit under the same gag without the visible footprint.

\textbf{H1b (differential timing).} The uncertainty component should follow the information clock, not the transaction clock: announcement ends the information asymmetry, so the component should resolve there; the cash position serves the purchase itself, so it should persist to completion.

These are predictions about where and when uncertain language appears, derived from the disclosure state; none asserts a cause, and we claim no identification.

\subsection{Estimation of the Main Variable}\label{subsec:mainvar}

Our main variable, UncResCEO, is the call-level residual of the Dzielinski--Wagner--Zeckhauser decomposition of CEO question-and-answer uncertainty: the component that remains after removing the CEO's persistent style (their equation 4). Only this unpersistent, call-varying shock in uncertainty language is usable for our question: an anticipatory, deal-tracking component must come and go within a CEO's tenure, so it can appear only in the residual. UncPreCEO, the raw uncertainty-word share of the prepared presentation, serves as the within-call control --- exactly its role in their equation 4: prepared remarks are drafted and vetted in advance, so language that responds to live questioning should load on the Q\&A side, not the script. Whether this residual tracks real uncertainty at all is an empirical question we answer before using it; Section~\ref{subsec:keyconstructs} shows it moves with established uncertainty measures and finds no evidence it is an artifact of analyst scrutiny.

\subsection{Methodology and Empirical Design}

Our designs stay close to the raw comparison. The run-up test is a within-firm comparison: in a two-way fixed-effects panel, we ask whether a cash acquirer's residual Q\&A uncertainty is higher in the quarter immediately preceding its announcement than in its own other quarters, with never-acquirers as the fixed-effects baseline --- a pre-window mean shift, not a pre/post difference-in-differences. The stock arm repeats the design for stock-financed acquirers, whose deals impose the same disclosure constraint without the visible cash footprint. The timing test is a disclosure-window event study on a matched universe: quarters are classified by the deal's life --- pre-announcement, the announced-but-pending gap, post-completion --- and the path of residual uncertainty is read against the path of the cash ratio over the same windows. Cash-specificity is tested formally: a pooled model interacts the pre-announcement indicator with the financing form and tests the cash--stock difference directly (Wald; Table~\ref{tab:empire_cashspec}). Fixed effects, clustering, and sample filters are reported with each table.

\subsection{Specification and Measurement of Key Constructs}\label{subsec:keyconstructs}

We ask two things of the residual: that it move with real uncertainty, and that it not be an artifact of analyst scrutiny.

The first is convergent validity. Firm-level political risk \citep{hassan2020} and aggregate policy uncertainty \citep{baker2016,davis2016} load positively on the residual (Tables~\ref{tab:h11_prisk_uncertainty}, \ref{tab:h24_us_epu}, and~\ref{tab:h24b_global_epu}) --- the same forces that move established uncertainty measures move this one. But these forces move the prepared presentation as well, so co-movement alone does not show the residual is distinct. The decisive evidence is discriminant: product-market competition \citep{hoberg2010,hoberg2016} loads on the presentation and not on the residual (Table~\ref{tab:h23_competition_uncertainty}). A persistent industry condition surfaces in the script, while the residual stays reserved for the call-varying shock.

The second is the analyst-scrutiny alternative: that pre-deal quarters simply draw harder questions about cash, and hard questions produce hedged answers. Section~\ref{sec:additional} tests it in three steps --- analysts do scrutinize cash (Table~\ref{tab:cash_scrutiny_validity}); scrutiny alone does not move the residual (Table~\ref{tab:cash_scrutiny_channel}); and the operative test, interacting scrutiny with the pre-announcement window, leaves the run-up intact (Table~\ref{tab:reason_gating}). We read this as a failure to find paired with a within-model dissociation, not a powered equivalence test.

\section{Main Empirical Analyses}\label{sec:main}

\subsection{Data, Sample, and Variable Construction}

Our data combine four sources. The textual core is the Capital IQ earnings-call transcript archive, covering the calls of US public firms from 2002 to 2018. Each transcript is parsed at the level of the individual speaking turn: every turn carries the speaker's role --- the CEO, other executives, or the questioning analysts --- and a segment marker separating the prepared presentation from the question-and-answer session, the split that allows scripted and unscripted CEO language to enter the analysis as distinct variables.

Uncertainty language is measured with the ``uncertainty'' category of the Loughran--McDonald Master Dictionary \citep{lm2011}; we use the 1993--2024 release. Firm fundamentals come from quarterly Compustat North America files, matched to each call as the most recent fiscal quarter ending on or before the call date; daily CRSP data and IBES earnings records supply the stock-return, market-return, and earnings-surprise controls that the decomposition in Section~\ref{subsec:mainvar} requires. Acquisition records come from SDC: for every deal we observe the announcement date, the completion or withdrawal date, and the payment composition that classifies it as cash- or stock-financed. SDC identifies acquirers by CUSIP while the call panel is keyed by Compustat's firm identifier, so the two are linked through the acquirer's six-digit CUSIP via the CRSP--Compustat link table.

We build the uncertainty measures ourselves. Each call is tokenized and the uncertainty-word share computed separately by speaker and segment; the CEO's question-and-answer share, UncAnsCEO, is the input to the main variable. Following \citet{dwz}, we regress UncAnsCEO on CEO fixed effects, the call's other speech characteristics (the presentation share UncPreCEO, analyst-question uncertainty, and negative tone), firm controls (stock and market returns, earnings growth, and the earnings surprise), and year effects --- their equation 4. The CEO fixed effect absorbs the executive's persistent style, and its negative is the clarity measure; UncResCEO is the residual, the call-specific innovation in unscripted CEO uncertainty that neither style nor circumstances explain. UncPreCEO, the raw uncertainty-word share of the prepared presentation (their equation 1), is retained as the prepared-remarks measure.

PreAnnounceQtr equals one in the single quarter immediately before the firm's first qualifying deal announcement ($e=-1$) and zero otherwise; never-acquirers serve as the fixed-effects baseline. CashRatio is cash and equivalents over total assets. CashScrutiny is the share of a call's analyst question turns containing at least one cash-level or liquidity term from a locked lexicon reproduced in the tables section; HighCashScrutiny flags any such scrutiny, since the median call draws none. The remaining controls --- leverage, size, Tobin's Q, profitability, investment, payout, and cash-flow volatility --- are defined, with every other variable, in the Appendix.

The estimation samples are built in three layers. The base layer is the call panel: every Capital IQ call of a US public firm between 2002 and 2018 that can be matched to Compustat, excluding financial and utility firms throughout and excluding CEOs with fewer than five calls in the sample, for whom no style can be estimated and hence no residual is defined. The second layer is the deal universe: an SDC acquisition qualifies if it is announced by a US public acquirer between 2002 and 2018, its status is completed, pending, or withdrawn, and its payment composition is known; deals at least 50\% cash-financed define the cash arm and deals at least 50\% stock-financed the stock placebo arm, with each firm's first qualifying deal setting its event clock. The full eligibility criteria are reproduced on the deal-specification page of the tables section (the page preceding Table~\ref{tab:empire_building_did}). The third layer is design-specific: the run-up test drops a treated firm's post-announcement quarters, the event studies cap the post-announcement window at four quarters, and the scrutiny validity and channel tests require at least three analyst Q\&A turns per call. Sample sizes therefore vary across designs because each test keeps only the firms for which its comparison is defined; every table reports its own firm and observation counts. Table~\ref{tab:summary_stats} reports summary statistics for the main estimation samples.

\subsection{Main Analysis 1: The Pre-Announcement Run-Up}

Main Analysis 1 asks whether the gag window is visible at all (Table~\ref{tab:empire_building_did}). We estimate $Y_{i,t} = \beta\,\mathrm{PreAnnounceQtr}_{i,t} + \boldsymbol{\gamma}'\mathbf{X}_{i,t} + \alpha_i + \tau_t + \varepsilon_{i,t}$ by two-way fixed-effects OLS --- firm and calendar year-quarter effects, standard errors clustered by firm --- four times: each outcome (the cash ratio and the residual) in each arm (cash and stock acquirers), holding the equation fixed (the full layout is on the specification page preceding the table). Identification is within-firm: a treated firm's pre-announcement quarter is compared with its own ordinary quarters and with never-acquirers; its post-announcement quarters are dropped, so the design observes only the run-up. Because the hypothesis is directional, the treatment coefficient is tested one-tailed ($\beta>0$); controls are two-tailed.

The cash-ratio equation behaves like a standard cash regression before the treatment enters: more levered, larger, and capital-intensive firms hold less cash, growth opportunities raise it, and the partial-adjustment lag is 0.7990 (SE 0.0070) --- cash is sticky, and with the lag absorbing the level the equation fits tightly ($R^2$ 0.665). The residual equation explains almost nothing by construction ($R^2$ 0.001): its dependent variable is already an innovation, purged of CEO style and call circumstances, so no control should have much left to say --- and none does, every control being insignificant in the cash arm.

In the cash arm, both clocks move in the final pre-announcement quarter. The cash ratio is still rising --- 0.0051 (SE 0.0017, 1\% level), a coefficient that tracks the change in cash, not its level, because of the lag --- and residual Q\&A uncertainty is elevated: 0.0461 (SE 0.0172, 1\%). The two columns run on different universes --- 67,590 firm-quarters of 2,232 firms for cash, 27,622 of 1,248 for the residual --- because the residual exists only for CEOs with at least five calls, for whom style is estimable; the uncertainty test is correspondingly the more demanding one. The scrutiny columns of the same table are taken up in Section~\ref{sec:additional}.

The effects are material. The uncertainty coefficient equals fifteen percent of a standard deviation of the residual (0.3072, Table~\ref{tab:summary_stats}, Panel B) --- for a variable from which everything persistent has already been removed, in a quarter that makes up 1.4\% of the sample. The cash coefficient says the firm adds half a percentage point of assets to its cash stock in that single quarter, roughly three percent of the mean cash ratio (0.1573, Panel A). This run-up is the trace predicted by H1.

The stock arm provides the placebo: its CEOs sit under the same disclosure constraint, but the uncertainty indicator shows no comparable positive run-up ($-0.0429$, SE 0.0307, n.s.) and the cash ratio no comparable accumulation ($-0.0036$, n.s.). The arms are otherwise alike --- the stock arm's partial-adjustment lag, 0.8013, is nearly identical --- so what differs at $e=-1$ is the response, not the model. Deal pendency alone, without the cash footprint, leaves no trace. A side-by-side reading is not yet a formal test of the difference; Main Analysis 3 supplies that.
% numbers verified programmatically against thesis_tables.tex (tmp/verify_draft_numbers.py); scrutiny columns of the same table reserved for §IV

\subsection{Main Analysis 2: Differential Timing Around the Announcement}

Main Analysis 2 replaces the single pre-announcement indicator with an event-time decomposition, so the same finding reads as a sequence rather than a point (Table~\ref{tab:empire_drop_matched}). On the matched universe --- 28,102 firm-quarters of 1,320 firms, those for which both the residual and the cash ratio are defined --- each treated quarter falls into one of four bins relative to the firm's first deal: two quarters out ($t-2$), the pre-announcement quarter ($t-1$), the announced-but-unclosed window, and the post-completion window. The specification is otherwise that of Main Analysis 1, and with no ordering imposed across bins every coefficient is two-tailed.

Two features license the reading. The $t-2$ coefficient is null on both clocks --- 0.0068 for the residual, 0.0008 for the cash ratio, neither significant --- so nothing moves two quarters out and the $t-1$ estimates are not the tail of a pre-existing trend. The cash equation again carries its partial-adjustment lag (0.7547) and its controls load as before, leverage and size entering the cash ratio negatively, so the cash clock is identified off the same well-behaved equation as in Main Analysis 1.

The two clocks then diverge in timing. Residual uncertainty spikes at $t-1$ (0.0473, 1\% level) and collapses the instant the deal is public: the announced-but-unclosed coefficient is 0.0018 (n.s.), a fall of 0.0455 (5\%) from the pre-announcement quarter, and by completion the level sits below baseline ($-0.0250$, 10\%) for a peak-to-post swing of 0.0723 (1\%). The cash ratio moves on the transaction, not the news --- elevated at $t-1$ (0.0061, 5\%), it does not fall when the deal is announced (drop 0.0006, n.s.) and declines only once cash leaves the firm at completion (post-level $-0.0155$; announce-to-completion drop 0.0210, 1\%).

The entry spike is economically the same event as Main Analysis 1 --- 0.0473 against 0.0461 there, about fifteen percent of a residual standard deviation. What Main Analysis 2 adds is the round trip: the full peak-to-post swing (0.0723) is half again the size of that entry spike, so the uncertainty raised in the gag window is entirely unwound once the information is out.

The claim rests on the contrast between the two paths, not the magnitude of any single drop --- uncertainty tracks the information and reverses at announcement, cash tracks the transaction and falls at completion. This is the differential-timing prediction of H1b. A side-by-side reading of cash against stock is still not a formal test of their difference; Main Analysis 3 supplies it.
% numbers verified programmatically against thesis_tables.tex (tmp/verify_draft_numbers.py)

\subsection{Main Analysis 3: Cash-Specificity}

Main Analysis 3 asks whether the run-up is specific to cash, in two steps. The first reruns the event study of Main Analysis 2 side by side on the stock placebo arm (Table~\ref{tab:empire_drop_placebo}); the second pools the two arms in a single interacted model that formally tests the cash--stock difference, for the effect (pre-announcement uncertainty) and for its proposed cause (the cash build-up) alike (Table~\ref{tab:empire_cashspec}). Both tables are two-tailed throughout.

For the reason given in Main Analysis 1 --- stock acquirers carry the identical disclosure gag but no cash war-chest --- a cash-driven effect should surface in the cash arm and be absent from the stock arm. Neither arm trends into the event: the $t-2$ pre-trend is flat in both (0.0105 for cash, $-0.0056$ for stock, neither significant).

Side by side, only the cash arm moves. It reproduces the run-up --- pre-announcement uncertainty of 0.0486 (1\%) that collapses to 0.0058 (n.s.) once the deal is public, a peak-to-post drop of 0.0681 (1\%) across 29,535 firm-quarters --- while the stock arm shows none of it: its pre-announcement coefficient is negative ($-0.0404$, n.s.; 39,819 firm-quarters), and if anything its uncertainty moves the other way across announcement (pre-to-gap difference $-0.0756$, 10\%).

A side-by-side display, though, implies a comparison it never runs, so the pooled model tests it directly. On the effect, the cash arm's pre-announcement coefficient of 0.0459 (5\%) against $-0.0524$ (n.s.) for stock yields a formal cash--stock difference of 0.0983 (5\%; 25,600 firm-quarters); because the arms move in opposite directions the gap exceeds either coefficient, about a third of a residual standard deviation. On the proposed cause, the cash build-up is only weakly cash-sided --- the difference is insignificant on the matched universe (0.0064, n.s.) and reaches just 0.0092 (10\%) on the full panel.

One formal test at the 5\% level supports H1a's claim of concentration in cash; the cause's own cash-specificity is weaker still, marginal and only in the full panel. We therefore read the result as supported but fragile --- concentration, not strict specificity --- and our wording stays with ``concentrated.''
% numbers verified programmatically against thesis_tables.tex (tmp/verify_draft_numbers.py)

\section{Additional Analyses}\label{sec:additional}

\subsection{Ruling Out Analyst Scrutiny}

The scrutiny alternative holds that pre-deal quarters draw harder analyst questions about cash, and that hard questions, not the deal, produce the hedged answers. Three results address it in order: whether the scrutiny measure is real, whether scrutiny moves the residual at all, and whether it accounts for this particular run-up.

The measure is real (Table~\ref{tab:cash_scrutiny_validity}). The share of analyst Q\&A turns devoted to cash and liquidity rises strongly with the firm's cash ratio (0.7530, 1\%, one-tailed; 75,087 calls) and with a high-cash indicator (0.1754, 1\%), and the cash-ratio link survives the full control set (0.8519). Analysts do watch the war chest, so the nulls that follow are not the artifact of a dead measure.

Yet scrutiny does not move the residual (Table~\ref{tab:cash_scrutiny_channel}). Its direct effect on UncResCEO is zero to four decimal places ($-0.0000$, SE 0.0013; 41,512 calls), and the null holds under a logit on the dichotomized residual ($-0.0003$) and on the raw uncertainty measure and coarser scrutiny proxies alike.

The operative test puts both in the same model on the cash-acquirer pre-announcement universe (Table~\ref{tab:reason_gating}). The pre-announcement indicator stays at 0.0413 (1\%) with scrutiny entered as a regressor and 0.0439 (5\%) with the interaction added, while scrutiny's own effect is null and the scrutiny~$\times$~pre-announcement interaction is $-0.0056$ (SE 0.0111, n.s.). The reason raises uncertainty; the questioning does not, and the questioning does not amplify the reason.

We keep the register honest. This is a failure to find paired with a within-model dissociation, not proof of absence: the interaction's confidence interval ($[-0.027, +0.016]$) does not exclude effects that would matter against a run-up near 0.04, 89\% of calls draw no cash scrutiny at all, and a significant-versus-insignificant contrast is not itself a test. The claim is that scrutiny does not account for \emph{this} run-up --- not that scrutiny never matters anywhere.
% numbers verified programmatically against thesis_tables.tex (tmp/verify_draft_numbers.py)

\subsection{The Presentation-Side Contrast}

The presentation segment supplies the contrast. If the residual is the deal-tracking component, then the scripted presentation should carry the persistent signal that outsiders price and regulators notice, and the residual should not --- a split two outcomes put to the test: the bid-ask spread and the receipt of an SEC comment letter.

On the bid-ask spread (Table~\ref{tab:h14c_ceo2_decomp}), the prepared remarks' uncertainty is priced. UncPreCEO enters at 0.1664 (5\%) in the baseline and 0.3496 (1\%) with firm fixed effects, fading as the full control set and finer time effects enter (industry track 0.0108, n.s.; firm track 0.1644, 10\%) --- a pattern with established precedent \citep{bushee2018}. The residual never is: UncResCEO is insignificant in all twelve spread specifications (baseline $-0.0594$, SE 0.1068; 42,625 firm-quarters).

Regulatory attention splits the same way (Table~\ref{tab:h18_ceo2_decomp}). Receipt of an SEC conference-call comment letter \citep{lerman2026} is weakly associated with presentation uncertainty (up to 0.0016, 5\%, across specifications) and never with the residual ($-0.0003$, n.s.; 44,113 firm-quarters).

Two components, two audiences: the scripted segment carries the persistent signal the market prices and regulators notice, the residual the episodic, deal-tracking one that neither does.

This is the contrast promised in Section~\ref{subsec:keyconstructs}, and it is deliberately not a headline: the effects are weak or fading, reported as texture, not pillars --- the load-bearing claims rest on the main analyses, not here.
% numbers verified programmatically against thesis_tables.tex (tmp/verify_draft_numbers.py)

\section{Conclusion}\label{sec:conclusion}

\subsection{Summary of Findings}

We set out to describe how CEO earnings-call language uncertainty behaves around undisclosed acquisitions. The picture is coherent: the unscripted residual of a CEO's question-and-answer language carries an anticipatory signal keyed to the disclosure state of a pending cash deal. It is present while the deal is still secret and concentrated among cash acquirers; it resolves the moment the deal becomes public, even as the acquirer's cash position keeps building toward completion; it is not an echo of analysts pressing harder on cash; and it sits in the call-varying residual, not in the scripted presentation that the market and regulators actually price. These are facets of one claim --- that call language tracks what a firm knows but has not yet said --- rather than five separate results.

We hold that claim at the strength the evidence supports: concentration in cash, not strict specificity, and correlation, not cause. The through-line is consistent across the run-up, its timing, and its resolution, even as the magnitudes stay modest and the formal test of cash-concentration rests on a single comparison. What the paper establishes is a pattern and its disclosure-state signature; the reading of that signature it leaves deliberately hedged.

\subsection{Contributions}

The paper's first contribution is a dimension and a place to look. The work nearest to it reads how corporate language around deals is \emph{managed} --- the tone of prepared remarks, the volume of strategy vocabulary; this paper reads what unmanaged language \emph{discloses}: the uncertainty that surfaces in unscripted answers when a CEO is bound to silence. Finding it required a specific place --- the call-varying residual, the only component in which an anticipatory, deal-tracking signal can survive once a CEO's persistent style and the firm's circumstances are stripped away. The cell it fills, cash acquirers and Q\&A uncertainty in the anticipatory window, is the one that work had left empty.

Its broader contribution is a lens. By starting from the firm's disclosure state rather than its balance sheet, the paper treats earnings-call language as carrying information about material events not yet disclosed, with a timing signature that appears under the gag and resolves at announcement --- language as a marker of what is known but not yet sayable, a reading more general than the acquisition setting that demonstrates it. What the paper offers is therefore a measurement and a pattern, not a mechanism: a dimension the field can now read and a state it can now track, with the reason the dimension moves --- compliance-constrained disclosure or strategic silence --- left deliberately open.

\subsection{Limitations}

Beyond the correlational stance, the unresolved mechanism, and the single formal test behind the cash concentration --- all acknowledged in their place --- two further limits bound what the paper can claim. The first is measurement. The signal is a language proxy: anticipatory uncertainty is inferred from the word-shares of unscripted answers, an indirect read of a private state the data never observe directly. The construct is validated --- it moves with established uncertainty measures and survives the scrutiny rule-out --- so the concern is not whether it captures uncertainty, but that it captures it at one remove, through language rather than through the knowledge or intent behind it.

The second is scope. The estimates come from US public acquirers between 2002 and 2018, observed on Capital IQ earnings calls, and from CEOs with enough calls for a persistent style to be estimated --- the only sample in which a call-varying residual is even defined. Private acquirers, other markets and periods, less-covered executives, and settings without a comparable disclosure regime lie outside it. What the paper documents, it documents within a particular institutional and linguistic frame.

\subsection{Directions for Future Research}

The clearest next step is a design that separates the two readings the present data cannot. Because compliance-constrained disclosure and strategic silence are observationally equivalent here, the lever is variation in the constraint itself: settings where the legal gag is tighter or looser, or where a CEO's private incentive to stay silent varies independently of it, would make the two accounts predict different things. Variation of that kind would do double duty --- naming the mechanism and, where it is exogenous, moving the claim from correlation toward cause.

The lens is also not limited to acquisitions. Any material event a firm is committed to but cannot yet disclose --- pending litigation, an unannounced leadership succession --- places a CEO in the same bind, and the residual could be turned on those windows to ask whether anticipatory call-uncertainty is a general feature of undisclosed information rather than an artifact of deals. Tested across settings and disclosure regimes, that question would also map how far the pattern travels. What the paper can already say is narrower, and we think it is enough: in the unscripted minutes of a routine earnings call, the language of a CEO bound to silence already carries the trace of what cannot yet be said.

% References — every entry verbatim-verified against the paper's own first/title page
% (local PDFs in docs/papers/ + NotebookLM scoped quotes, 2026-06-09).
\begin{thebibliography}{10}

\bibitem[Baker et~al.(2016)]{baker2016}
Baker, S.~R., N.~Bloom, and S.~J. Davis. 2016. Measuring economic policy uncertainty. \emph{The Quarterly Journal of Economics} 131: 1593--1636.

\bibitem[Basic Inc. v.\ Levinson(1988)]{basic1988}
\emph{Basic Inc.\ v.\ Levinson}, 485 U.S.\ 224 (1988).

\bibitem[Bushee et~al.(2018)]{bushee2018}
Bushee, B.~J., I.~D. Gow, and D.~J. Taylor. 2018. Linguistic complexity in firm disclosures: Obfuscation or information? \emph{Journal of Accounting Research} 56: 85--121.

\bibitem[Davis(2016)]{davis2016}
Davis, S.~J. 2016. An index of global economic policy uncertainty. NBER Working Paper 22740.

\bibitem[Dzielinski et~al.(2021)]{dwz}
Dzielinski, M., A.~F. Wagner, and R.~J. Zeckhauser. 2021. Straight talkers and vague talkers: The effects of managerial style in earnings conference calls. M-RCBG Faculty Working Paper Series 2017-02, Harvard Kennedy School (2017, revised 2021).

\bibitem[Everhart et~al.(2025)]{everhart2025}
Everhart, J., T.~Kravet, S.~McVay, and J.~Warren. 2025. The impact of M\&A transactions on acquiring firm guidance. Working paper, August 2025.

\bibitem[Gokkaya et~al.(2025)]{gokkaya2025}
Gokkaya, S., X.~Liu, and R.~M. Stulz. 2025. Is there information in corporate acquisition plans? Charles A. Dice Center Working Paper 2024-04, June 2025.

\bibitem[Hassan et~al.(2020)]{hassan2020}
Hassan, T.~A., S.~Hollander, L.~van Lent, and A.~Tahoun. 2020. Firm-level political risk: Measurement and effects. \emph{The Quarterly Journal of Economics} 135: 2135--2202.

\bibitem[Hoberg and Phillips(2010)]{hoberg2010}
Hoberg, G., and G.~Phillips. 2010. Product market synergies and competition in mergers and acquisitions: A text-based analysis. \emph{Review of Financial Studies} 23: 3773--3811.

\bibitem[Hoberg and Phillips(2016)]{hoberg2016}
Hoberg, G., and G.~Phillips. 2016. Text-based network industries and endogenous product differentiation. \emph{Journal of Political Economy} 124: 1423--1465.

\bibitem[Lerman et~al.(2026)]{lerman2026}
Lerman, A., T.~D. Steffen, and K.~Zhang. 2026. Earnings conference calls and the SEC comment letter process. \emph{Management Science}, Articles in Advance: 1--34.

\bibitem[Loughran and McDonald(2011)]{lm2011}
Loughran, T., and B.~McDonald. 2011. When is a liability not a liability? Textual analysis, dictionaries, and 10-Ks. \emph{The Journal of Finance} 66: 35--65.

\bibitem[Matsumoto et~al.(2011)]{matsumoto2011}
Matsumoto, D., M.~Pronk, and E.~Roelofsen. 2011. What makes conference calls useful? The information content of managers' presentations and analysts' discussion sessions. \emph{The Accounting Review} 86: 1383--1414.

\bibitem[Ragozzino and Reuer(2024)]{ragozzino2024}
Ragozzino, R., and J.~J. Reuer. 2024. Implications of mergers and acquisitions for information disclosures in earnings calls. \emph{Long Range Planning} 57: 102393.

\bibitem[Rule 10b-5(2014)]{rule10b5}
U.S. Securities and Exchange Commission. Employment of manipulative and deceptive devices. Rule 10b-5, 17 C.F.R.\ \S\,240.10b-5.

\bibitem[Thewissen et~al.(2024)]{thewissen2024}
Thewissen, J., B.~Yan, \"O.~Arslan-Ayaydin, and S.~Yan. 2024. Manipulating disclosure tone: Understanding acquiring firms' strategies in stock-for-stock mergers and acquisitions. SSRN Working Paper 4900453.

\bibitem[Bertrand and Schoar(2003)]{bertrand_schoar2003}
Bertrand, M., and A.~Schoar. 2003. Managing with style: The effect of managers on firm policies. \emph{The Quarterly Journal of Economics} 118: 1169--1208.

\bibitem[Dye(1985)]{dye1985}
Dye, R. 1985. Disclosure of nonproprietary information. \emph{Journal of Accounting Research} 23: 123--145.

\bibitem[Harford(1999)]{harford1999}
Harford, J. 1999. Corporate cash reserves and acquisitions. \emph{The Journal of Finance} 54: 1969--1997.

\bibitem[Hollander et~al.(2010)]{hollander2010}
Hollander, S., M.~Pronk, and E.~Roelofsen. 2010. Does silence speak? An empirical analysis of disclosure choices during conference calls. \emph{Journal of Accounting Research} 48: 531--563.

\bibitem[Keown and Pinkerton(1981)]{keown1981}
Keown and Pinkerton. 1981. Merger announcements and insider trading activity: An empirical investigation. \emph{The Journal of Finance} 36: 855--869.

\bibitem[Verrecchia(1983)]{verrecchia1983}
Verrecchia, R. 1983. Discretionary disclosure. \emph{Journal of Accounting and Economics} 5: 179--194.

\end{thebibliography}

\appendix
\section{Variable Definitions and Measurement}
% Definitions copied from the bible's specification pages (_empire_building_spec.tex +
% the _cash_scrutiny_validity.tex construction page) and the DWZ source paper — 2026-06-09.

\noindent Uncertainty words are the ``uncertainty'' category of the Loughran--McDonald Master Dictionary \citep{lm2011}; we use the 1993--2024 release.

\vspace{4pt}
\begin{center}
\small
\begin{tabular}{p{0.27\columnwidth}p{0.62\columnwidth}}
\toprule
Variable & Definition \\
\midrule
UncResCEO & residual from the \citet{dwz} equation-4 regression of the CEO's Q\&A uncertainty-word share on CEO fixed effects, speech and firm controls, and year effects --- the call-varying component net of the CEO's persistent style \\
UncPreCEO & raw uncertainty-word share of the CEO's presentation segment (their equation 1) \\
CashRatio & $\mathrm{cheq}/\mathrm{atq}$ \\
PreAnnounceQtr & $=1$ in the single quarter immediately before the firm's first qualifying deal ($e=-1$); post-announcement quarters dropped; never-acquirers always $0$ \\
CashScrutiny & share (\%) of a call's analyst Q\&A turns containing at least one cash-level or liquidity term \\
HighCash & $\mathbf{1}[\mathrm{CashRatio} \geq \text{top tercile}]$ \\
Leverage & $(\mathrm{dlcq}+\mathrm{dlttq})/\mathrm{atq}$ \\
lnAssets & $\ln(\mathrm{atq})$ \\
TobinsQ & $(\mathrm{atq}+\mathrm{cshoq}\cdot\mathrm{prccq}-\mathrm{ceqq})/\mathrm{atq}$ \\
ROA & $\mathrm{iby}/\overline{\mathrm{assets}}$ \\
Capex & $\mathrm{capxy}/\mathrm{atq}$ \\
DivDummy & $\mathbb{1}[\mathrm{dvy}>0]$ \\
sCFO & 5-yr rolling SD of $\mathrm{oancfy}/\mathrm{atq}_{t-1}$ \\
\midrule
\multicolumn{2}{l}{\textit{Event-window and treatment variables}} \\
PRE2, PRE1, GAP, POST & event-window indicators: $e=-2$; $e=-1$; announced but not yet closed; completed (call quarter $\geq$ the deal's effective-date quarter); post-announcement window capped at $+4$ quarters \\
PreAnn\_cash, PreAnn\_stock & pooled-model arm indicators: $=1$ in the quarter before the firm's first $\geq$50\%-cash ($\geq$50\%-stock) deal \\
CashRatio$_{t-1}$ & within-firm one-quarter lag of CashRatio (consecutive quarters only); partial-adjustment control in the CashRatio regressions \\
\midrule
\multicolumn{2}{l}{\textit{Language variables}} \\
UncAnsCEO & raw uncertainty-word share of the CEO's Q\&A turns (the equation-4 dependent variable) \\
ClarityCEO & negative of the CEO fixed effect from equation 4 \\
HighCashScrutiny & $\mathbf{1}[\mathrm{CashScrutiny} > \text{sample median}]$; the median is zero, so it flags any cash scrutiny \\
UncQue & uncertainty-word share of the analyst Q\&A turns \\
NegCall & negative-word share of the entire call \\
\midrule
\multicolumn{2}{l}{\textit{Validity-test variables (Section~\ref{sec:additional})}} \\
PRisk & firm-level political risk (Hassan et al.), matched by firm and calendar quarter \\
US\_EPU\_log & $\ln$ of the Baker--Bloom--Davis US economic policy uncertainty index \\
GEPU\_log & $\ln$ of the Davis global economic policy uncertainty index \\
z\_log\_TotalSimilarity & z-scored $\ln$ of Hoberg--Phillips TNIC3 total product-market similarity \\
\midrule
\multicolumn{2}{l}{\textit{Market-quality and regulatory variables (Section~\ref{sec:additional})}} \\
BGTLevel\_Spread & mean daily relative spread $2(\mathrm{ASK}-\mathrm{BID})/(\mathrm{ASK}+\mathrm{BID})$ over the call day and the following 25 trading days, $\times 10^{4}$ \\
CCCL & $\mathbf{1}$[SEC conference-call comment letter received between the current and the next call] \\
Lagged\_DV & the regression's own dependent variable measured at the prior quarter \\
\midrule
\multicolumn{2}{l}{\textit{Additional controls (Section~\ref{sec:additional})}} \\
FirmMat & $\mathrm{req}/\mathrm{atq}$ \\
EarnVol & SD of annual ROA over the trailing five fiscal years \\
StockPrice & CRSP closing price at the call date (nearest prior trading day) \\
Turnover & daily volume over shares outstanding at the call date \\
DailyVola & annualized SD (\%) of daily returns over the window between consecutive calls \\
AbsSurpDec & absolute value of SurpDec, the within-quarter $-5$ to $+5$ rank of the IBES actual-minus-consensus surprise \\
SalesGrowth & $(\mathrm{sale}_{t}-\mathrm{sale}_{t-1})/|\mathrm{sale}_{t-1}|$, annual ($\mathrm{sale}=\mathrm{saley}$) \\
RDSales & $\mathrm{xrdy}/\mathrm{saley}$ \\
CashFlowAt & $\mathrm{oancfy}/\overline{\mathrm{assets}}$ \\
\bottomrule
\end{tabular}
\end{center}

\noindent Full construction detail --- the complete search-term lists and deal-eligibility criteria --- is reported on the specification pages of the tables document.

\clearpage
% Tables: docs/Draft/thesis_tables.tex is the ONE table/number source. The blocks below are
% BYTE-EXACT copies generated by tmp/extract_draft_tables.py — regenerate, never hand-edit.
\section*{Tables}
% AUTO-GENERATED by tmp/extract_draft_tables.py — BYTE-EXACT copies from the bible
% (docs/Draft/thesis_tables.tex + fragments). DO NOT HAND-EDIT. Regenerate instead.

% --- tab:summary_stats
\begin{table}[htbp]
\centering
\caption{Summary Statistics}
\label{tab:summary_stats}
\scriptsize
\begin{tabular}{lrrrrrr}
\toprule
 & N & Mean & SD & P25 & Median & P75 \\
\midrule
\multicolumn{7}{l}{\textit{Panel A. Cash-acquirer universe, CashRatio equation (run-up test, column 1)}} \\
\addlinespace
CashRatio & 67,590 & 0.1573 & 0.1834 & 0.0270 & 0.0843 & 0.2200 \\
CashRatio$_{t-1}$ & 67,590 & 0.1578 & 0.1842 & 0.0270 & 0.0842 & 0.2207 \\
PreAnnounceQtr & 67,590 & 0.0108 & 0.1035 & 0.0000 & 0.0000 & 0.0000 \\
Leverage & 67,590 & 0.2421 & 0.2377 & 0.0491 & 0.2062 & 0.3685 \\
lnAssets & 67,590 & 7.5823 & 1.7055 & 6.3609 & 7.5367 & 8.7235 \\
TobinsQ & 67,590 & 2.0287 & 1.5770 & 1.1486 & 1.5053 & 2.2687 \\
ROA & 67,590 & 0.0356 & 0.1296 & 0.0100 & 0.0412 & 0.0859 \\
Capex & 67,590 & 0.0465 & 0.0558 & 0.0114 & 0.0305 & 0.0613 \\
DivDummy & 67,590 & 0.5551 & 0.4970 & 0.0000 & 1.0000 & 1.0000 \\
sCFO & 67,590 & 0.0638 & 0.2866 & 0.0184 & 0.0357 & 0.0649 \\
\midrule
\multicolumn{7}{l}{\textit{Panel B. Cash-acquirer universe, UncResCEO equation (run-up test, column 2)}} \\
\addlinespace
UncResCEO & 27,622 & 0.0002 & 0.3072 & -0.1869 & -0.0174 & 0.1686 \\
CashScrutiny & 26,216 & 0.3867 & 1.4701 & 0.0000 & 0.0000 & 0.0000 \\
HighCashScrutiny & 26,216 & 0.1115 & 0.3147 & 0.0000 & 0.0000 & 0.0000 \\
PreAnnounceQtr & 27,622 & 0.0143 & 0.1189 & 0.0000 & 0.0000 & 0.0000 \\
Leverage & 27,622 & 0.2220 & 0.2346 & 0.0283 & 0.1919 & 0.3345 \\
lnAssets & 27,622 & 7.2945 & 1.4587 & 6.2459 & 7.2395 & 8.2685 \\
TobinsQ & 27,622 & 2.2300 & 1.5195 & 1.3143 & 1.7408 & 2.5956 \\
ROA & 27,622 & 0.0539 & 0.1141 & 0.0266 & 0.0612 & 0.1037 \\
Capex & 27,622 & 0.0535 & 0.0557 & 0.0199 & 0.0362 & 0.0647 \\
DivDummy & 27,622 & 0.5028 & 0.5000 & 0.0000 & 1.0000 & 1.0000 \\
sCFO & 27,622 & 0.0662 & 0.2524 & 0.0247 & 0.0406 & 0.0674 \\
\bottomrule
\end{tabular}
\begin{minipage}{\linewidth}
\vspace{2pt}\scriptsize
\textit{Notes:} Moments are computed on the exact estimation samples of the Empire-Building Run-Up Test (cash-acquirer arm): Panel A on the CashRatio-equation universe (column 1), Panel B on the UncResCEO-equation universe (column 2); the CashScrutiny and HighCashScrutiny rows use their matched universe (columns 3--4), hence their smaller N. Never-acquirer firm-quarters are included (the fixed-effects baseline); treated firms' post-announcement quarters are excluded by design. All variables are defined in the Appendix.
\end{minipage}
\end{table}
\clearpage

% --- deal-specification page (precedes the run-up table, as cited in the text)
% Hand-written specification page for the Empire-Building Run-Up Test.
% Static explanatory content (NOT generated). Placed before _empire_building_did.tex.
\begin{center}
{\large\bfseries Empire-Building Run-Up Test --- Regression Specification}
\end{center}

\vspace{4pt}
\noindent All four regressions share one estimating equation; they differ only in the
outcome $Y$ and in which firms are eligible to be treated. Estimator: two-way
fixed-effects OLS (\texttt{PanelOLS}), firm-clustered standard errors.

\begin{equation*}
Y_{i,t} \;=\; \beta\,\mathrm{PreAnnounceQtr}_{i,t}
\;+\; \boldsymbol{\gamma}'\mathbf{X}_{i,t}
\;+\; \alpha_i \;+\; \tau_t \;+\; \varepsilon_{i,t}
\end{equation*}

\noindent where $i$ indexes the firm (its earnings call) and $t$ the calendar
year-quarter. $\beta$ is the coefficient of interest: one-tailed test ($H_1:\beta>0$)
for $\mathrm{PreAnnounceQtr}$, two-tailed for the controls.

\begin{center}
\small
\begin{tabular}{cll}
\toprule
Reg. & Outcome $Y$ & Treated set (firms eligible for $\mathrm{PreAnnounceQtr}=1$) \\
\midrule
(1) & CashRatio  & Cash acquirers ($\geq 50\%$ cash) \\
(2) & UncResCEO  & Cash acquirers ($\geq 50\%$ cash) \\
(3) & CashRatio  & Stock acquirers ($\geq 50\%$ stock) --- placebo \\
(4) & UncResCEO  & Stock acquirers ($\geq 50\%$ stock) --- placebo \\
\bottomrule
\end{tabular}
\end{center}

\begin{center}
\small
\begin{tabular}{lll}
\toprule
Symbol & Definition & Role \\
\midrule
\multicolumn{3}{l}{\textit{Outcomes}} \\
CashRatio & $\mathrm{cheq}/\mathrm{atq}$ (Compustat) & cash-holdings level (first stage) \\
UncResCEO & residual $\widehat\varepsilon_{ic}$ from $\mathrm{UncAnsCEO}_{ic}=\alpha+\mathrm{CEO\text{-}FE}_i+\mathbf{X}_{ic}\boldsymbol\gamma+\varepsilon_{ic}$ (DWZ Eqn.\ 4) & call-level CEO Q\&A uncertainty, net of the CEO's permanent style \\
\midrule
\multicolumn{3}{l}{\textit{Treatment}} \\
$\mathrm{PreAnnounceQtr}_{i,t}$ & $=1$ in the single quarter immediately before firm $i$'s first qualifying deal ($e=-1$); & pre-event (anticipation) indicator \\
 & post-announcement quarters ($e\geq 0$) are dropped; never-acquirers are always $0$ & \\
\midrule
\multicolumn{3}{l}{\textit{Controls $\mathbf{X}$ (H1 cash set, lagged DV removed; Compustat quarterly)}} \\
Leverage & $(\mathrm{dlcq}+\mathrm{dlttq})/\mathrm{atq}$ & interest-bearing debt / assets \\
lnAssets & $\ln(\mathrm{atq})$ & firm size \\
TobinsQ  & $(\mathrm{atq}+\mathrm{cshoq}\cdot\mathrm{prccq}-\mathrm{ceqq})/\mathrm{atq}$ & market-to-book / growth opportunities \\
ROA      & $\mathrm{iby}/\overline{\mathrm{assets}}$ & profitability \\
Capex    & $\mathrm{capxy}/\mathrm{atq}$ & investment intensity \\
DivDummy & $\mathbb{1}[\mathrm{dvy}>0]$ & dividend payer \\
sCFO     & 5-yr rolling SD of $\mathrm{oancfy}/\mathrm{atq}_{t-1}$ & cash-flow volatility \\
\midrule
\multicolumn{3}{l}{\textit{Fixed effects / errors}} \\
$\alpha_i$ & firm fixed effect (\texttt{EntityEffects}) & removes each firm's permanent level (within-firm identification) \\
$\tau_t$   & calendar year-quarter fixed effect (\texttt{TimeEffects}) & removes economy-wide shocks each quarter \\
$\varepsilon_{i,t}$ & idiosyncratic error, clustered by firm & \\
\bottomrule
\end{tabular}
\end{center}

\vspace{2pt}
\noindent\small\textit{Reading the columns.} (1)--(2) are the cash-acquirer arm: does cash, and
then CEO uncertainty, rise in the quarter before a cash-financed deal? (3)--(4) are the
stock-acquirer placebo: stock buyers face the same pending-deal disclosure gag but hold
no cash war-chest, so (4) flat while (2) positive $\Rightarrow$ the uncertainty rise is
cash-specific, not mere deal secrecy.
\clearpage

% --- tab:empire_building_did
\begin{table}[htbp]
\centering
\caption{Empire-Building Run-Up Test}
\label{tab:empire_building_did}
\scriptsize
\begin{tabular}{lcccccccc}
\toprule
 & \multicolumn{4}{c}{Cash acquirers} & \multicolumn{4}{c}{Stock acquirers (placebo)} \\
\cmidrule(lr){2-5}\cmidrule(lr){6-9}
 & (1) & (2) & (3) & (4) & (5) & (6) & (7) & (8) \\
 & CashRatio & UncResCEO & CashScrutiny & HighCashScrutiny & CashRatio & UncResCEO & CashScrutiny & HighCashScrutiny \\
\midrule
PreAnnounceQtr & \textbf{0.0051}$^{***}$ & \textbf{0.0461}$^{***}$ & 0.0854 & \textbf{0.0408}$^{**}$ & -0.0036 & -0.0429 & -0.3222 & -0.0303 \\
 & (0.0017) & (0.0172) & (0.0749) & (0.0194) & (0.0035) & (0.0307) & (0.2024) & (0.0313) \\
CashRatio$_{t-1}$ (partial adj.) & \textbf{0.7990}$^{***}$ & --- & --- & --- & \textbf{0.8013}$^{***}$ & --- & --- & --- \\
 & (0.0070) &  &  &  & (0.0057) &  &  &  \\
\midrule
Leverage & \textbf{-0.0146}$^{***}$ & -0.0338 & -0.1521 & -0.0179 & \textbf{-0.0213}$^{***}$ & -0.0292 & -0.1056 & -0.0428 \\
 & (0.0039) & (0.0217) & (0.1470) & (0.0308) & (0.0034) & (0.0196) & (0.1358) & (0.0262) \\
lnAssets & \textbf{-0.0072}$^{***}$ & -0.0037 & -0.0391 & \textbf{0.0190}$^{**}$ & \textbf{-0.0079}$^{***}$ & -0.0043 & \textbf{-0.0733}$^{**}$ & 0.0077 \\
 & (0.0011) & (0.0082) & (0.0482) & (0.0096) & (0.0008) & (0.0063) & (0.0369) & (0.0078) \\
TobinsQ & \textbf{0.0013}$^{**}$ & -0.0031 & -0.0208 & 0.0006 & \textbf{0.0015}$^{***}$ & -0.0017 & -0.0113 & 0.0030 \\
 & (0.0006) & (0.0030) & (0.0144) & (0.0027) & (0.0004) & (0.0025) & (0.0123) & (0.0024) \\
ROA & 0.0053 & 0.0499 & \textbf{-0.5659}$^{**}$ & \textbf{-0.1006}$^{**}$ & \textbf{0.0075}$^{*}$ & \textbf{0.0530}$^{*}$ & \textbf{-0.5350}$^{***}$ & \textbf{-0.0869}$^{**}$ \\
 & (0.0042) & (0.0321) & (0.2254) & (0.0405) & (0.0041) & (0.0296) & (0.1982) & (0.0366) \\
Capex & \textbf{-0.1057}$^{***}$ & 0.0842 & \textbf{-1.2215}$^{***}$ & \textbf{-0.3646}$^{***}$ & \textbf{-0.1067}$^{***}$ & 0.1256 & \textbf{-1.4479}$^{***}$ & \textbf{-0.4470}$^{***}$ \\
 & (0.0126) & (0.0986) & (0.3534) & (0.0880) & (0.0113) & (0.1006) & (0.3000) & (0.0793) \\
DivDummy & -0.0011 & -0.0030 & -0.0561 & -0.0111 & \textbf{-0.0017}$^{**}$ & \textbf{-0.0137}$^{*}$ & -0.0564 & \textbf{-0.0146}$^{*}$ \\
 & (0.0009) & (0.0106) & (0.0537) & (0.0107) & (0.0008) & (0.0080) & (0.0384) & (0.0087) \\
sCFO & 0.0008 & 0.0052 & \textbf{0.1180}$^{***}$ & \textbf{0.0235}$^{***}$ & 0.0007 & -0.0009 & \textbf{0.0906}$^{***}$ & \textbf{0.0169}$^{***}$ \\
 & (0.0006) & (0.0085) & (0.0140) & (0.0028) & (0.0007) & (0.0110) & (0.0337) & (0.0057) \\
\midrule
Firm FE & Yes & Yes & Yes & Yes & Yes & Yes & Yes & Yes \\
Cal. Year-Quarter FE & Yes & Yes & Yes & Yes & Yes & Yes & Yes & Yes \\
\midrule
Firms & 2,232 & 1,248 & 1,237 & 1,237 & 2,281 & 1,338 & 1,325 & 1,325 \\
N (firm-quarters) & 67,590 & 27,622 & 26,216 & 26,216 & 84,979 & 39,377 & 37,065 & 37,065 \\
$R^2$ & 0.665 & 0.001 & 0.002 & 0.002 & 0.672 & 0.001 & 0.002 & 0.002 \\
\bottomrule
\end{tabular}
\begin{minipage}{\linewidth}
\vspace{2pt}\scriptsize
\textit{Notes:} $^{*}p<0.10$, $^{**}p<0.05$, $^{***}p<0.01$ (one-tailed for the treatment coefficient, $\beta > 0$; two-tailed for controls). 
Significant coefficients in \textbf{bold}. Standard errors (in parentheses) clustered at firm level. CashScrutiny (cols 3, 7) and HighCashScrutiny (cols 4, 8) are estimated on the UncResCEO call universe. HighCashScrutiny is an LPM on $\mathbf{1}$[CashScrutiny $>$ median]; the median is $0$ (89\% of calls raise no cash turns), so it flags any cash scrutiny ($\approx$11\% of calls). The CashRatio columns (1, 5) add CashRatio's own one-quarter lag (partial-adjustment for the sticky DV), so the PreAnnounceQtr coefficient there tracks the \emph{change} in cash, not its level; the other DVs take no lag.
\end{minipage}
\end{table}
\clearpage

% --- tab:empire_drop_matched
\begin{table}[htbp]
\centering
\caption{Pre-Acquisition Uncertainty vs.\ Cash on the Matched Universe (disclosure-window event study)}
\label{tab:empire_drop_matched}
\scriptsize
\begin{tabular}{lcc}
\toprule
 & (1) UncResCEO & (2) CashRatio \\
\midrule
PRE2 ($t{-}2$, pre-trend) & 0.0068 & 0.0008 \\
 & (0.0178) & (0.0024) \\
PRE1 ($t{-}1$, pre-announce) & \textbf{0.0473}$^{***}$ & \textbf{0.0061}$^{**}$ \\
 & (0.0178) & (0.0024) \\
GAP (announced, pre-close) & 0.0018 & 0.0055 \\
 & (0.0187) & (0.0034) \\
POST (completed) & \textbf{-0.0250}$^{*}$ & \textbf{-0.0155}$^{***}$ \\
 & (0.0130) & (0.0022) \\
\midrule
Drop: PRE1 $-$ GAP & \textbf{0.0455}$^{**}$ & 0.0006 \\
 & (0.0225) & (0.0039) \\
Drop: GAP $-$ POST & 0.0268 & \textbf{0.0210}$^{***}$ \\
 & (0.0199) & (0.0042) \\
Drop: PRE1 $-$ POST & \textbf{0.0723}$^{***}$ & \textbf{0.0216}$^{***}$ \\
 & (0.0190) & (0.0031) \\
\midrule
\multicolumn{3}{l}{\textit{Controls}} \\
Leverage & \textbf{-0.0415}$^{*}$ & \textbf{-0.0215}$^{***}$ \\
 & (0.0213) & (0.0054) \\
lnAssets & -0.0060 & \textbf{-0.0074}$^{***}$ \\
 & (0.0084) & (0.0020) \\
TobinsQ & -0.0045 & \textbf{0.0013}$^{**}$ \\
 & (0.0029) & (0.0006) \\
ROA & \textbf{0.0572}$^{*}$ & 0.0067 \\
 & (0.0325) & (0.0061) \\
Capex & 0.1073 & \textbf{-0.1434}$^{***}$ \\
 & (0.1054) & (0.0141) \\
DivDummy & -0.0099 & -0.0016 \\
 & (0.0104) & (0.0014) \\
sCFO & \textbf{0.0096}$^{**}$ & -0.0011 \\
 & (0.0043) & (0.0007) \\
CashRatio$_{t-1}$ (partial adj.) & --- & \textbf{0.7547}$^{***}$ \\
 &  & (0.0108) \\
\midrule
Firm FE / Year-Qtr FE / Controls & Yes & Yes \\
N (firm-quarters) & 28,102 & 28,102 \\
Firms & 1,320 & 1,320 \\
\bottomrule
\end{tabular}
\begin{minipage}{\linewidth}\vspace{2pt}\scriptsize
\textit{Notes:} $^{*}p<.10$, $^{**}p<.05$, $^{***}p<.01$ (two-tailed).
\end{minipage}
\end{table}
\clearpage

% --- tab:empire_cashspec
\begin{table}[htbp]
\centering
\caption{Formal Cash-Specificity: Pre-Announcement Uncertainty (effect) vs.\ Cash Build-Up (proposed cause), cash vs.\ stock acquirers}
\label{tab:empire_cashspec}
\scriptsize
\begin{tabular}{lccc}
\toprule
 & \multicolumn{1}{c}{EFFECT} & \multicolumn{2}{c}{PROPOSED CAUSE} \\
\cmidrule(lr){2-2}\cmidrule(lr){3-4}
 & (1) UncResCEO & (2) CashRatio & (3) CashRatio \\
 & (matched) & (matched, +lag) & (full panel, +lag) \\
\midrule
Pre-announce qtr, Cash acquirer & \textbf{0.0459}$^{**}$ & \textbf{0.0056}$^{**}$ & \textbf{0.0049}$^{***}$ \\
 & (0.0185) & (0.0025) & (0.0019) \\
Pre-announce qtr, Stock acquirer & -0.0524 & -0.0008 & -0.0043 \\
 & (0.0436) & (0.0066) & (0.0051) \\
\midrule
Cash $-$ Stock (formal test) & \textbf{0.0983}$^{**}$ & 0.0064 & \textbf{0.0092}$^{*}$ \\
 & (0.0476) & (0.0071) & (0.0055) \\
\midrule
\multicolumn{4}{l}{\textit{Controls}} \\
Leverage & \textbf{-0.0411}$^{*}$ & \textbf{-0.0236}$^{***}$ & \textbf{-0.0194}$^{***}$ \\
 & (0.0230) & (0.0056) & (0.0043) \\
lnAssets & -0.0034 & -0.0033 & \textbf{-0.0065}$^{***}$ \\
 & (0.0092) & (0.0021) & (0.0011) \\
TobinsQ & -0.0025 & 0.0007 & \textbf{0.0012}$^{**}$ \\
 & (0.0031) & (0.0006) & (0.0005) \\
ROA & 0.0556 & 0.0046 & 0.0054 \\
 & (0.0355) & (0.0065) & (0.0046) \\
Capex & 0.0487 & \textbf{-0.1437}$^{***}$ & \textbf{-0.1095}$^{***}$ \\
 & (0.1065) & (0.0148) & (0.0139) \\
DivDummy & -0.0068 & \textbf{-0.0026}$^{*}$ & -0.0015 \\
 & (0.0109) & (0.0014) & (0.0010) \\
sCFO & 0.0040 & -0.0008 & 0.0009 \\
 & (0.0088) & (0.0007) & (0.0006) \\
\quad CashRatio$_{t-1}$ (partial adj.) & --- & \textbf{0.7778}$^{***}$ & \textbf{0.7957}$^{***}$ \\
 &  & (0.0113) & (0.0073) \\
\midrule
Firm FE / Year-Qtr FE / Controls & Yes & Yes & Yes \\
UncResCEO-present restriction & Yes & Yes & \textbf{No} \\
N (firm-quarters) & 25,600 & 24,347 & 61,876 \\
Firms & 1,191 & 1,176 & 2,141 \\
\bottomrule
\end{tabular}
\begin{minipage}{\linewidth}\vspace{2pt}\scriptsize
\textit{Notes:} $^{*}p<.10$, $^{**}p<.05$, $^{***}p<.01$ (two-tailed).
\end{minipage}
\end{table}
\clearpage

% --- tab:reason_gating
\begin{table}[htbp]
\centering
\caption{Cash-Scrutiny Reason-Gating Test}
\label{tab:reason_gating}
\small
\begin{tabular}{lcc}
\toprule
 & \multicolumn{2}{c}{UncResCEO} \\
\cmidrule(lr){2-3}
 & (1) & (2) \\
\midrule
CashScrutiny & 0.0000 & 0.0001 \\
 & (0.0018) & (0.0018) \\
PreAnnounceQtr & \textbf{0.0413}$^{***}$ & \textbf{0.0439}$^{**}$ \\
 & (0.0177) & (0.0189) \\
CashScrutiny $\times$ PreAnnounceQtr &  & -0.0056 \\
 &  & (0.0111) \\
\midrule
Controls & Yes & Yes \\
Firm FE & Yes & Yes \\
Cal. Year-Quarter FE & Yes & Yes \\
\midrule
N (firm-quarters) & 26,216 & 26,216 \\
Firms & 1,237 & 1,237 \\
$R^2$ & 0.001 & 0.001 \\
\bottomrule
\end{tabular}
\begin{minipage}{\linewidth}
\vspace{2pt}\scriptsize
\textit{Notes:} $^{*}p<0.10$, $^{**}p<0.05$, $^{***}p<0.01$ (one-tailed, $\beta>0$, for the directional terms); significant coefficients in \textbf{bold}. Standard errors (in parentheses) clustered at firm level. Estimated on cash-acquirer firms ($\geq$50\%-cash acquisitions; SDC, US public, 2002--2018) over their pre-announcement and earlier quarters (post-announcement quarters dropped); \texttt{PreAnnounceQtr} $=$ $\mathbf{1}$[the single quarter before announcement]. Sample restricted to calls with both \texttt{CashScrutiny} and \texttt{UncResCEO} (the matched universe). \texttt{CashScrutiny} $=$ \% of a call's analyst Q\&A turns on cash/liquidity. The reason (\texttt{PreAnnounceQtr}) raises CEO residual uncertainty; analyst cash-scrutiny does not, and the interaction is insignificant.
\end{minipage}
\end{table}
\clearpage

% --- tab:h11_prisk_uncertainty
\begin{table}[htbp]
\centering
\caption{Political Risk and Call Uncertainty}
\label{tab:h11_prisk_uncertainty}
\scriptsize
\begin{tabular}{lcccc}
\toprule
 & (1) & (2) & (3) & (4) \\
 & \multicolumn{2}{c}{Industry FE} & \multicolumn{2}{c}{Firm FE} \\
\cmidrule(lr){2-3} \cmidrule(lr){4-5}
 & UncResCEO & UncPreCEO & UncResCEO & UncPreCEO \\
\midrule
PRisk & \textbf{0.0001}$^{***}$ & \textbf{0.0003}$^{***}$ & \textbf{0.0001}$^{***}$ & \textbf{0.0002}$^{***}$ \\
 & (0.0000) & (0.0000) & (0.0000) & (0.0000) \\
\midrule
UncQue & -0.0022 & \textbf{0.0204}$^{***}$ & -0.0014 & \textbf{0.0081}$^{***}$ \\
 & (0.0032) & (0.0045) & (0.0037) & (0.0029) \\
NegCall & -0.0025 & \textbf{0.1874}$^{***}$ & -0.0036 & \textbf{0.1477}$^{***}$ \\
 & (0.0051) & (0.0141) & (0.0073) & (0.0082) \\
lnAssets & -0.0002 & \textbf{-0.0296}$^{***}$ & -0.0054 & 0.0048 \\
 & (0.0004) & (0.0039) & (0.0058) & (0.0078) \\
TobinsQ & -0.0005 & -0.0007 & -0.0008 & 0.0001 \\
 & (0.0009) & (0.0036) & (0.0022) & (0.0023) \\
ROA & \textbf{0.0291}$^{**}$ & 0.0477 & \textbf{0.0540}$^{**}$ & 0.0346 \\
 & (0.0145) & (0.0318) & (0.0264) & (0.0230) \\
CashRatio & 0.0019 & -0.0354 & 0.0206 & 0.0265 \\
 & (0.0077) & (0.0357) & (0.0240) & (0.0287) \\
DivDummy & -0.0031 & -0.0139 & -0.0108 & \textbf{-0.0229}$^{**}$ \\
 & (0.0021) & (0.0118) & (0.0074) & (0.0103) \\
FirmMat & 0.0012 & \textbf{0.0021}$^{**}$ & 0.0027 & 0.0006 \\
 & (0.0013) & (0.0010) & (0.0041) & (0.0007) \\
EarnVol & 0.0078 & -0.0072 & 0.0056 & -0.0104 \\
 & (0.0077) & (0.0157) & (0.0117) & (0.0132) \\
UncPreCEO & -0.0051 &  & -0.0054 &  \\
 & (0.0034) &  & (0.0054) &  \\
\midrule
Industry FE & Yes & Yes &  &  \\
Firm FE &  &  & Yes & Yes \\
Year FE & Yes & Yes & Yes & Yes \\
\midrule
N & 44,497 & 65,760 & 44,497 & 65,760 \\
$R^2$ & 0.003 & 0.043 & 0.003 & 0.023 \\
Adj.~$R^2$ & 0.002 & 0.043 & -0.030 & -0.004 \\
\bottomrule
\end{tabular}
\begin{minipage}{\linewidth}
\vspace{2pt}\scriptsize
\textit{Notes:} $^{*}p<0.10$, $^{**}p<0.05$, $^{***}p<0.01$ (one-tailed for IVs, $\beta > 0$; two-tailed for controls).
 Significant coefficients in \textbf{bold}.
 Standard errors (in parentheses) clustered at firm level.
 Main sample (excludes financial and utility firms). Firms with fewer than 5 calls are excluded.
 $R^2$ includes absorbed fixed effects (not within-$R^2$).
\end{minipage}
\end{table}
\clearpage

% --- tab:h24_us_epu
\begin{table}[htbp]
\centering
\caption{US Policy Uncertainty and Call Uncertainty}
\label{tab:h24_us_epu}
\scriptsize
\begin{tabular}{lcccc}
\toprule
 & (1) & (2) & (3) & (4) \\
 & \multicolumn{2}{c}{Industry + Cal. Year FE} & \multicolumn{2}{c}{Firm + Cal. Year FE} \\
\cmidrule(lr){2-3} \cmidrule(lr){4-5}
 & UncResCEO & UncPreCEO & UncResCEO & UncPreCEO \\
\midrule
US\_EPU\_log & \textbf{0.0124}$^{**}$ & \textbf{0.0211}$^{***}$ & \textbf{0.0123}$^{*}$ & \textbf{0.0239}$^{***}$ \\
 & (0.0075) & (0.0074) & (0.0075) & (0.0076) \\
\midrule
UncQue & 0.0011 & \textbf{0.0164}$^{***}$ & 0.0024 & \textbf{0.0095}$^{***}$ \\
 & (0.0035) & (0.0032) & (0.0040) & (0.0028) \\
NegCall & 0.0006 & \textbf{0.1247}$^{***}$ & 0.0021 & \textbf{0.1385}$^{***}$ \\
 & (0.0056) & (0.0089) & (0.0073) & (0.0079) \\
lnAssets & 0.0009 & \textbf{-0.0143}$^{***}$ & 0.0013 & 0.0038 \\
 & (0.0007) & (0.0020) & (0.0067) & (0.0063) \\
TobinsQ & -0.0001 & 0.0007 & -0.0000 & 0.0007 \\
 & (0.0011) & (0.0019) & (0.0027) & (0.0018) \\
ROA & \textbf{0.0281}$^{*}$ & \textbf{0.0367}$^{**}$ & \textbf{0.0677}$^{**}$ & \textbf{0.0467}$^{**}$ \\
 & (0.0149) & (0.0168) & (0.0268) & (0.0191) \\
CashRatio & 0.0051 & -0.0079 & 0.0345 & 0.0305 \\
 & (0.0087) & (0.0187) & (0.0256) & (0.0231) \\
DivDummy & \textbf{-0.0046}$^{***}$ & -0.0074 & \textbf{-0.0142}$^{**}$ & \textbf{-0.0157}$^{*}$ \\
 & (0.0009) & (0.0061) & (0.0072) & (0.0083) \\
FirmMat & \textbf{0.0023}$^{*}$ & 0.0009 & 0.0005 & -0.0011 \\
 & (0.0012) & (0.0008) & (0.0047) & (0.0012) \\
EarnVol & \textbf{0.0268}$^{*}$ & -0.0074 & \textbf{0.0353}$^{**}$ & -0.0083 \\
 & (0.0157) & (0.0086) & (0.0176) & (0.0119) \\
UncPreCEO & -0.0026 &  & -0.0011 &  \\
 & (0.0032) &  & (0.0045) &  \\
Lagged\_DV & \textbf{0.0202}$^{***}$ & \textbf{0.5176}$^{***}$ & \textbf{0.0187}$^{***}$ & \textbf{0.2626}$^{***}$ \\
 & (0.0054) & (0.0131) & (0.0056) & (0.0089) \\
\midrule
Industry FE & Yes & Yes &  &  \\
Firm FE &  &  & Yes & Yes \\
Year FE & Yes & Yes & Yes & Yes \\
\midrule
N & 39,063 & 60,503 & 39,063 & 60,503 \\
$R^2$ & 0.001 & 0.292 & 0.001 & 0.084 \\
Adj.~$R^2$ & -0.000 & 0.291 & -0.035 & 0.057 \\
\bottomrule
\end{tabular}
\begin{minipage}{\linewidth}
\vspace{2pt}\scriptsize
\textit{Notes:} $^{*}p<0.10$, $^{**}p<0.05$, $^{***}p<0.01$ (one-tailed for IVs, $\beta > 0$; two-tailed for controls).
 Significant coefficients in \textbf{bold}.
 Standard errors (in parentheses) two-way clustered (firm, calendar quarter).
 Main sample (excludes financial and utility firms). Firms with fewer than 5 calls are excluded.
 $R^2$ includes absorbed fixed effects (not within-$R^2$).
\end{minipage}
\end{table}
\clearpage

% --- tab:h24b_global_epu
\begin{table}[htbp]
\centering
\caption{Global Policy Uncertainty and Call Uncertainty}
\label{tab:h24b_global_epu}
\scriptsize
\begin{tabular}{lcccc}
\toprule
 & (1) & (2) & (3) & (4) \\
 & \multicolumn{2}{c}{Industry + Cal. Year FE} & \multicolumn{2}{c}{Firm + Cal. Year FE} \\
\cmidrule(lr){2-3} \cmidrule(lr){4-5}
 & UncResCEO & UncPreCEO & UncResCEO & UncPreCEO \\
\midrule
GEPU\_log & \textbf{0.0181}$^{**}$ & \textbf{0.0271}$^{***}$ & \textbf{0.0187}$^{**}$ & \textbf{0.0309}$^{***}$ \\
 & (0.0100) & (0.0107) & (0.0101) & (0.0108) \\
\midrule
UncQue & 0.0011 & \textbf{0.0164}$^{***}$ & 0.0024 & \textbf{0.0095}$^{***}$ \\
 & (0.0035) & (0.0032) & (0.0039) & (0.0028) \\
NegCall & 0.0006 & \textbf{0.1247}$^{***}$ & 0.0020 & \textbf{0.1385}$^{***}$ \\
 & (0.0056) & (0.0089) & (0.0072) & (0.0079) \\
lnAssets & 0.0009 & \textbf{-0.0143}$^{***}$ & 0.0012 & 0.0037 \\
 & (0.0007) & (0.0020) & (0.0067) & (0.0063) \\
TobinsQ & -0.0001 & 0.0008 & 0.0000 & 0.0007 \\
 & (0.0011) & (0.0019) & (0.0027) & (0.0018) \\
ROA & \textbf{0.0280}$^{*}$ & \textbf{0.0367}$^{**}$ & \textbf{0.0673}$^{**}$ & \textbf{0.0462}$^{**}$ \\
 & (0.0149) & (0.0168) & (0.0267) & (0.0191) \\
CashRatio & 0.0051 & -0.0080 & 0.0345 & 0.0304 \\
 & (0.0087) & (0.0188) & (0.0256) & (0.0231) \\
DivDummy & \textbf{-0.0047}$^{***}$ & -0.0074 & \textbf{-0.0142}$^{**}$ & \textbf{-0.0156}$^{*}$ \\
 & (0.0009) & (0.0061) & (0.0072) & (0.0083) \\
FirmMat & \textbf{0.0023}$^{*}$ & 0.0009 & 0.0006 & -0.0011 \\
 & (0.0012) & (0.0008) & (0.0047) & (0.0012) \\
EarnVol & \textbf{0.0269}$^{*}$ & -0.0074 & \textbf{0.0354}$^{**}$ & -0.0082 \\
 & (0.0157) & (0.0086) & (0.0176) & (0.0119) \\
UncPreCEO & -0.0026 &  & -0.0011 &  \\
 & (0.0032) &  & (0.0045) &  \\
Lagged\_DV & \textbf{0.0202}$^{***}$ & \textbf{0.5176}$^{***}$ & \textbf{0.0187}$^{***}$ & \textbf{0.2626}$^{***}$ \\
 & (0.0054) & (0.0131) & (0.0056) & (0.0089) \\
\midrule
Industry FE & Yes & Yes &  &  \\
Firm FE &  &  & Yes & Yes \\
Year FE & Yes & Yes & Yes & Yes \\
\midrule
N & 39,063 & 60,503 & 39,063 & 60,503 \\
$R^2$ & 0.001 & 0.292 & 0.001 & 0.084 \\
Adj.~$R^2$ & -0.000 & 0.291 & -0.035 & 0.057 \\
\bottomrule
\end{tabular}
\begin{minipage}{\linewidth}
\vspace{2pt}\scriptsize
\textit{Notes:} $^{*}p<0.10$, $^{**}p<0.05$, $^{***}p<0.01$ (one-tailed for IVs, $\beta > 0$; two-tailed for controls).
 Significant coefficients in \textbf{bold}.
 Standard errors (in parentheses) two-way clustered (firm, calendar quarter).
 Main sample (excludes financial and utility firms). Firms with fewer than 5 calls are excluded.
 $R^2$ includes absorbed fixed effects (not within-$R^2$).
\end{minipage}
\end{table}
\clearpage

% --- tab:h23_competition_uncertainty
\begin{table}[htbp]
\centering
\caption{Product-Market Competition and Call Uncertainty}
\label{tab:h23_competition_uncertainty}
\scriptsize
\begin{tabular}{lcccc}
\toprule
 & (1) & (2) & (3) & (4) \\
 & \multicolumn{2}{c}{Industry FE} & \multicolumn{2}{c}{Firm FE} \\
\cmidrule(lr){2-3} \cmidrule(lr){4-5}
 & UncResCEO & UncPreCEO & UncResCEO & UncPreCEO \\
\midrule
z\_log\_TotalSimilarity & 0.0008 & \textbf{0.0304}$^{***}$ & 0.0023 & \textbf{0.0302}$^{***}$ \\
 & (0.0014) & (0.0074) & (0.0075) & (0.0093) \\
\midrule
UncQue & 0.0011 & \textbf{0.0398}$^{***}$ & 0.0039 & 0.0084 \\
 & (0.0062) & (0.0100) & (0.0087) & (0.0073) \\
NegCall & 0.0020 & \textbf{0.2340}$^{***}$ & 0.0046 & \textbf{0.1698}$^{***}$ \\
 & (0.0068) & (0.0201) & (0.0120) & (0.0140) \\
lnAssets & -0.0007 & \textbf{-0.0335}$^{***}$ & -0.0101 & 0.0013 \\
 & (0.0007) & (0.0038) & (0.0071) & (0.0081) \\
TobinsQ & 0.0001 & -0.0007 & 0.0009 & -0.0017 \\
 & (0.0013) & (0.0037) & (0.0025) & (0.0025) \\
ROA & 0.0233 & \textbf{0.0681}$^{**}$ & 0.0343 & 0.0241 \\
 & (0.0197) & (0.0314) & (0.0308) & (0.0254) \\
CashRatio & 0.0065 & \textbf{-0.0803}$^{**}$ & 0.0055 & 0.0022 \\
 & (0.0103) & (0.0360) & (0.0287) & (0.0306) \\
DivDummy & -0.0020 & -0.0060 & -0.0074 & -0.0166 \\
 & (0.0028) & (0.0117) & (0.0083) & (0.0106) \\
FirmMat & 0.0018 & \textbf{0.0014}$^{**}$ & 0.0083 & 0.0003 \\
 & (0.0019) & (0.0007) & (0.0056) & (0.0004) \\
EarnVol & -0.0047 & -0.0210 & -0.0048 & -0.0114 \\
 & (0.0087) & (0.0152) & (0.0121) & (0.0146) \\
UncPreCEO & \textbf{0.0111}$^{**}$ &  & \textbf{0.0230}$^{**}$ &  \\
 & (0.0044) &  & (0.0093) &  \\
\midrule
Industry FE & Yes & Yes &  &  \\
Firm FE &  &  & Yes & Yes \\
Year FE & Yes & Yes & Yes & Yes \\
\midrule
N & 12,728 & 18,492 & 12,728 & 18,492 \\
$R^2$ & 0.001 & 0.050 & 0.001 & 0.021 \\
Adj.~$R^2$ & -0.002 & 0.048 & -0.098 & -0.068 \\
\bottomrule
\end{tabular}
\begin{minipage}{\linewidth}
\vspace{2pt}\scriptsize
\textit{Notes:} $^{*}p<0.10$, $^{**}p<0.05$, $^{***}p<0.01$ (one-tailed for IVs, $\beta > 0$; two-tailed for controls).
 Significant coefficients in \textbf{bold}.
 Standard errors (in parentheses) clustered at firm level.
 Main sample (excludes financial and utility firms). Unit of observation: firm-fiscal-year.
 $R^2$ includes absorbed fixed effects (not within-$R^2$).
\end{minipage}
\end{table}
\clearpage

% --- tab:cash_scrutiny_validity
\begin{table}[htbp]
\centering
\caption{Analyst Cash-Scrutiny Validity (Link 1)}
\label{tab:cash_scrutiny_validity}
\scriptsize
\begin{tabular}{lcccc}
\toprule
& \multicolumn{4}{c}{DV: CashScrutiny (\% of analyst Q\&A turns on cash/liquidity)} \\
\cmidrule(lr){2-5}
 & (1) & (2) & (3) & (4) \\
\midrule
CashRatio & \textbf{0.7530}$^{***}$ & \textbf{0.8519}$^{***}$ &  &  \\
 & (0.0949) & (0.0967) &  &  \\
HighCash &  &  & \textbf{0.1754}$^{***}$ & \textbf{0.1921}$^{***}$ \\
 &  &  & (0.0237) & (0.0236) \\
\midrule
Leverage & --- & \textbf{0.2180}$^{**}$ & --- & \textbf{0.2030}$^{**}$ \\
 &  & (0.0978) &  & (0.0987) \\
lnAssets & --- & \textbf{-0.0562}$^{***}$ & --- & \textbf{-0.0744}$^{***}$ \\
 &  & (0.0213) &  & (0.0212) \\
TobinsQ & --- & \textbf{-0.0272}$^{***}$ & --- & \textbf{-0.0218}$^{**}$ \\
 &  & (0.0089) &  & (0.0085) \\
ROA & --- & \textbf{-0.5012}$^{***}$ & --- & \textbf{-0.5002}$^{***}$ \\
 &  & (0.1339) &  & (0.1335) \\
Capex & --- & \textbf{-1.2348}$^{***}$ & --- & \textbf{-1.3259}$^{***}$ \\
 &  & (0.2890) &  & (0.2900) \\
DivDummy & --- & \textbf{-0.0674}$^{**}$ & --- & \textbf{-0.0678}$^{**}$ \\
 &  & (0.0331) &  & (0.0331) \\
sCFO & --- & 0.0244 & --- & 0.0291 \\
 &  & (0.0274) &  & (0.0277) \\
\midrule
Firm FE & Yes & Yes & Yes & Yes \\
Cal. Year-Quarter FE & Yes & Yes & Yes & Yes \\
Controls & No & Yes & No & Yes \\
\midrule
Firms & 1,868 & 1,853 & 1,868 & 1,853 \\
N (calls) & 75,087 & 72,670 & 75,087 & 72,670 \\
$R^2$ & 0.002 & 0.005 & 0.001 & 0.004 \\
\bottomrule
\end{tabular}
\begin{minipage}{\linewidth}
\vspace{2pt}\scriptsize
\textit{Notes:} $^{*}p<0.10$, $^{**}p<0.05$, $^{***}p<0.01$ (one-tailed for the cash coefficient, $\beta>0$; two-tailed for controls).
Significant coefficients in \textbf{bold}.
Standard errors (in parentheses) clustered at firm level.
\end{minipage}
\end{table}
\clearpage
\begin{center}\large\textbf{Cash-Scrutiny Measure: Variable Construction (Link 1)}\end{center}
\vspace{4pt}\noindent\textbf{Variable definitions.}
\begin{itemize}\setlength\itemsep{2pt}
\item \textbf{CashScrutiny}$_{i,t}$ --- share (\%) of call $i$'s analyst Q\&A turns whose text contains at least one \emph{cash level} or \emph{liquidity} term (listed below); a turn-level binary averaged over the call's analyst Q\&A turns.
\item \textbf{CashRatio}$_{i,t}$ = cheq/atq (cash \& equivalents $\div$ total assets).
\item \textbf{High Cash}$_{i,t}$ = $\mathbf{1}[\text{CashRatio}_{i,t}\geq$ top tercile$]$.
\item \textbf{Controls} $X_{i,t}$: Leverage, ln(Assets), Tobin's Q, ROA, Capex, Dividend Payer, Cash Flow (sCFO).
\end{itemize}
\vspace{2pt}\noindent\textbf{Specification.} Two-way fixed-effects OLS estimated at the call level:
\[ \text{CashScrutiny}_{i,t} = \beta\,\text{Cash}_{i,t} + \gamma' X_{i,t} + \alpha_i + \delta_{q(t)} + \varepsilon_{i,t}, \]
\noindent where $\text{Cash}_{i,t}$ = CashRatio (cols 1--2) or High Cash (cols 3--4); $\alpha_i$ = firm fixed effects, $\delta_{q(t)}$ = calendar year-quarter fixed effects; standard errors clustered by firm. The cash coefficient $\beta$ is the external-validity test.
\vspace{6pt}\noindent\textbf{Search terms.} An analyst Q\&A turn is flagged if, after the \emph{excluded} phrases are blanked out, its lower-cased text contains any phrase below, matched on a word boundary (case-insensitive). Every literal spelling actually searched is listed: where a phrase has a singular/plural or hyphenated/closed variant (e.g.\ \texttt{'noncash'} and \texttt{'non-cash'}), both forms appear. Disposition/payout terms (dividends, buybacks) are deliberately \emph{not} used here: they load on firm size and maturity and would correlate with CashRatio spuriously rather than through cash attention.
\begin{itemize}\setlength\itemsep{3pt}
\item \textbf{Cash level} (22 forms) --- \emph{why:} they name the firm's cash stockpile or balance itself, the direct object of attention to how much cash the firm holds. \\ \texttt{'cash holding'}, \texttt{'cash holdings'}, \texttt{'cash balance'}, \texttt{'cash balances'}, \texttt{'cash position'}, \texttt{'cash on hand'}, \texttt{'cash on the balance sheet'}, \texttt{'cash reserve'}, \texttt{'cash reserves'}, \texttt{'cash and cash equivalents'}, \texttt{'cash and equivalents'}, \texttt{'cash and shortterm investments'}, \texttt{'cash and short-term investments'}, \texttt{'net cash'}, \texttt{'cash pile'}, \texttt{'cash hoard'}, \texttt{'cash stockpile'}, \texttt{'war chest'}, \texttt{'dry powder'}, \texttt{'excess cash'}, \texttt{'idle cash'}, \texttt{'surplus cash'}
\item \textbf{Liquidity} (5 forms) --- \emph{why:} standard analyst synonyms for the same liquid-asset buffer. \\ \texttt{'liquidity'}, \texttt{'liquid assets'}, \texttt{'shortterm investments'}, \texttt{'short-term investments'}, \texttt{'marketable securities'}
\item \textbf{Excluded} (17 forms) --- \emph{why:} they contain ``cash'' but denote cash \emph{flow}, earnings, an accounting basis, or an idiom rather than the cash holdings, so they are blanked out before matching to prevent false positives. \\ \texttt{'free cash flow'}, \texttt{'operating cash flow'}, \texttt{'cash flow statement'}, \texttt{'cash flow'}, \texttt{'cash flows'}, \texttt{'cash conversion cycle'}, \texttt{'cash conversion'}, \texttt{'cash basis'}, \texttt{'cash cow'}, \texttt{'cash compensation'}, \texttt{'noncash'}, \texttt{'non-cash'}, \texttt{'cash register'}, \texttt{'cash crop'}, \texttt{'cash taxes'}, \texttt{'cash earnings'}, \texttt{'cash in on'}
\end{itemize}
\clearpage

% --- tab:cash_scrutiny_channel
\begin{table}[htbp]
\centering
\caption{Analyst Cash-Scrutiny Channel (Link 2)}
\label{tab:cash_scrutiny_channel}
\small
\begin{tabular}{lcc}
\toprule
 & UncResCEO & UncAnsCEO \\
 & (DWZ residual) & (raw Q\&A) \\
\midrule
\multicolumn{3}{l}{\textit{Panel A. OLS --- DV in continuous level}} \\
CashScrutiny & -0.0000 & 0.0001 \\
 & (0.0013) & (0.0013) \\
$N$ (calls) & 41,512 & 41,512 \\
Within $R^2$ & 0.000 & 0.000 \\
\midrule
\multicolumn{3}{l}{\textit{Panel B. Logit --- DV $= \mathbf{1}$[measure $\geq$ median]}} \\
CashScrutiny (log-odds) & -0.0003 & 0.0080 \\
 & (0.0074) & (0.0085) \\
$N$ (calls) & 41,512 & 41,512 \\
Pseudo $R^2$ & 0.002 & 0.019 \\
\midrule
Firm FE (A) / Industry FE (B) & Yes & Yes \\
Cal. Year-Quarter FE & Yes & Yes \\
\bottomrule
\end{tabular}
\begin{minipage}{\linewidth}\vspace{2pt}\scriptsize
\textit{Notes:} DV is CEO Q\&A uncertainty as the DWZ residual (\texttt{UncResCEO}, firm/linguistic component removed) or the raw measure (\texttt{UncAnsCEO}, pooled-1\% winsorized). \texttt{CashScrutiny} $=$ the share (\%) of a call's Q\&A turns devoted to cash topics (validated in the Link~1 table). Panel~A is OLS on the continuous DV; Panel~B is a logit on $\mathbf{1}$[DV $\geq$ sample median]. $^{*}p<0.10$, $^{**}p<0.05$, $^{***}p<0.01$ (two-tailed); significant coefficients in \textbf{bold}; SE (parentheses) clustered at firm level. Coarser scrutiny measures (any-turn, raw count, log count) give the same null.
\end{minipage}
\end{table}
\clearpage

% --- tab:empire_drop_placebo
\begin{table}[htbp]
\centering
\caption{Pre-Acquisition UncResCEO Event Study: Cash vs.\ Stock Acquirers (placebo)}
\label{tab:empire_drop_placebo}
\scriptsize
\begin{tabular}{lcc}
\toprule
 & \multicolumn{2}{c}{DV: UncResCEO} \\
\cmidrule(lr){2-3}
 & (1) Cash acquirers & (2) Stock acquirers \\
\midrule
PRE2 ($t{-}2$, pre-trend) & 0.0105 & -0.0056 \\
 & (0.0174) & (0.0368) \\
PRE1 ($t{-}1$, pre-announce) & \textbf{0.0486}$^{***}$ & -0.0404 \\
 & (0.0174) & (0.0299) \\
GAP (announced, pre-close) & 0.0058 & 0.0353 \\
 & (0.0181) & (0.0318) \\
POST (completed) & -0.0195 & -0.0048 \\
 & (0.0127) & (0.0246) \\
\midrule
Drop: PRE1 $-$ GAP & \textbf{0.0428}$^{*}$ & \textbf{-0.0756}$^{*}$ \\
 & (0.0221) & (0.0405) \\
Drop: GAP $-$ POST & 0.0253 & 0.0401 \\
 & (0.0194) & (0.0382) \\
Drop: PRE1 $-$ POST & \textbf{0.0681}$^{***}$ & -0.0355 \\
 & (0.0187) & (0.0358) \\
\midrule
\multicolumn{3}{l}{\textit{Controls}} \\
Leverage & \textbf{-0.0372}$^{*}$ & -0.0286 \\
 & (0.0211) & (0.0194) \\
lnAssets & -0.0063 & -0.0041 \\
 & (0.0080) & (0.0062) \\
TobinsQ & -0.0034 & -0.0018 \\
 & (0.0029) & (0.0025) \\
ROA & \textbf{0.0534}$^{*}$ & \textbf{0.0543}$^{*}$ \\
 & (0.0315) & (0.0293) \\
Capex & 0.0816 & 0.1297 \\
 & (0.0993) & (0.0987) \\
DivDummy & -0.0065 & \textbf{-0.0132}$^{*}$ \\
 & (0.0100) & (0.0080) \\
sCFO & 0.0049 & -0.0008 \\
 & (0.0084) & (0.0110) \\
\midrule
Firm FE / Year-Qtr FE / Controls & Yes & Yes \\
N (firm-quarters) & 29,535 & 39,819 \\
Firms & 1,326 & 1,363 \\
\bottomrule
\end{tabular}
\begin{minipage}{\linewidth}\vspace{2pt}\scriptsize
\textit{Notes:} $^{*}p<.10$, $^{**}p<.05$, $^{***}p<.01$ (two-tailed).
\end{minipage}
\end{table}
\clearpage

% --- tab:h14c_ceo2_decomp
\begin{table}[htbp]
\centering
\caption{CEO Uncertainty and the Bid-Ask Spread}
\label{tab:h14c_ceo2_decomp}
\scriptsize
\begin{tabular}{lcccccccccccc}
\toprule
 & (1) & (2) & (3) & (4) & (5) & (6) & (7) & (8) & (9) & (10) & (11) & (12) \\
 & \multicolumn{6}{c}{BGTLevel\_Spread} & \multicolumn{6}{c}{BGTLevel\_Spread\_lead1} \\
\cmidrule(lr){2-7} \cmidrule(lr){8-13}
\midrule
ClarityCEO & -0.0477 & 1.0617 & 0.2122 & 1.1056 & 0.1644 & 1.0547 & -0.1886 & 0.5998 & -0.0922 & 0.6828 & -0.0431 & 0.5722 \\
 & (0.2145) & (0.4277) & (0.2227) & (0.4260) & (0.2103) & (0.4174) & (0.2370) & (0.4004) & (0.2399) & (0.4055) & (0.2206) & (0.3729) \\
UncResCEO & -0.0594 & -0.0687 & -0.0882 & -0.1021 & -0.0757 & -0.0893 & 0.0825 & 0.0696 & 0.0742 & 0.0545 & 0.0984 & 0.0751 \\
 & (0.1068) & (0.1057) & (0.1023) & (0.1007) & (0.0994) & (0.0981) & (0.1416) & (0.1437) & (0.1407) & (0.1426) & (0.1339) & (0.1357) \\
UncPreCEO & \textbf{0.1664}$^{**}$ & \textbf{0.3496}$^{***}$ & 0.0814 & \textbf{0.2945}$^{***}$ & 0.0108 & \textbf{0.1644}$^{*}$ & -0.0817 & -0.0288 & -0.1072 & -0.0360 & 0.0170 & 0.1124 \\
 & (0.0946) & (0.1253) & (0.0956) & (0.1229) & (0.0919) & (0.1184) & (0.0926) & (0.1229) & (0.0926) & (0.1233) & (0.0872) & (0.1186) \\
\midrule
lnAssets & \textbf{-0.8039}$^{***}$ & \textbf{-1.9098}$^{***}$ & \textbf{-0.6339}$^{***}$ & \textbf{-1.6825}$^{***}$ & \textbf{-0.6352}$^{***}$ & \textbf{-1.8277}$^{***}$ & \textbf{-0.8218}$^{***}$ & \textbf{-1.5191}$^{***}$ & \textbf{-0.7857}$^{***}$ & \textbf{-1.7292}$^{***}$ & \textbf{-0.6568}$^{***}$ & \textbf{-1.4181}$^{***}$ \\
 & (0.0490) & (0.2000) & (0.0472) & (0.2156) & (0.0473) & (0.2158) & (0.0513) & (0.1951) & (0.0505) & (0.2130) & (0.0474) & (0.2017) \\
TobinsQ & \textbf{-0.2861}$^{***}$ & \textbf{-0.4785}$^{***}$ & \textbf{-0.2633}$^{***}$ & \textbf{-0.5518}$^{***}$ & \textbf{-0.2654}$^{***}$ & \textbf{-0.5883}$^{***}$ & \textbf{-0.2305}$^{***}$ & \textbf{-0.3215}$^{***}$ & \textbf{-0.2285}$^{***}$ & \textbf{-0.4803}$^{***}$ & \textbf{-0.2035}$^{***}$ & \textbf{-0.4289}$^{***}$ \\
 & (0.0457) & (0.0900) & (0.0470) & (0.0963) & (0.0451) & (0.0956) & (0.0467) & (0.0850) & (0.0478) & (0.1002) & (0.0433) & (0.0935) \\
ROA & \textbf{-9.4469}$^{***}$ & \textbf{-11.6430}$^{***}$ & \textbf{-5.7059}$^{***}$ & \textbf{-7.6844}$^{***}$ & \textbf{-5.7511}$^{***}$ & \textbf{-7.8517}$^{***}$ & \textbf{-9.9256}$^{***}$ & \textbf{-12.8443}$^{***}$ & \textbf{-8.8744}$^{***}$ & \textbf{-11.9197}$^{***}$ & \textbf{-7.6375}$^{***}$ & \textbf{-10.5867}$^{***}$ \\
 & (1.0201) & (1.3867) & (0.9822) & (1.2654) & (0.9633) & (1.2490) & (1.0798) & (1.2992) & (1.0722) & (1.2732) & (0.9869) & (1.1921) \\
Leverage & \textbf{0.3875}$^{*}$ & \textbf{1.8330}$^{***}$ & -0.1126 & 0.6850 & -0.0542 & 0.8026 & 0.2924 & 0.5282 & 0.1338 & 0.2221 & 0.0592 & 0.1928 \\
 & (0.2088) & (0.5143) & (0.2434) & (0.5013) & (0.2264) & (0.4954) & (0.2224) & (0.5516) & (0.2248) & (0.5596) & (0.2111) & (0.5258) \\
Capex & -0.8001 & 0.9540 & \textbf{-2.3622}$^{**}$ & -1.5271 & \textbf{-2.0790}$^{*}$ & -1.2897 & \textbf{-1.6463}$^{*}$ & -3.2121 & \textbf{-2.1561}$^{**}$ & \textbf{-4.3379}$^{*}$ & -1.4770 & -2.4356 \\
 & (1.1109) & (2.9126) & (1.1044) & (2.8061) & (1.0845) & (2.7998) & (0.9620) & (2.2020) & (0.9823) & (2.2446) & (0.9011) & (2.1565) \\
DivDummy & 0.1043 & 0.1429 & \textbf{0.3329}$^{***}$ & 0.2549 & \textbf{0.3156}$^{***}$ & 0.2723 & 0.0240 & 0.1007 & 0.1012 & 0.1454 & 0.1145 & 0.1739 \\
 & (0.0893) & (0.2020) & (0.0896) & (0.1888) & (0.0855) & (0.1867) & (0.0887) & (0.2095) & (0.0889) & (0.2074) & (0.0825) & (0.1975) \\
sCFO & \textbf{0.6700}$^{**}$ & \textbf{0.8405}$^{**}$ & 0.4045 & \textbf{0.6825}$^{*}$ & 0.4084 & \textbf{0.6882}$^{**}$ & 0.0871 & \textbf{0.3644}$^{*}$ & 0.0040 & 0.3165 & -0.0840 & 0.2479 \\
 & (0.3405) & (0.3472) & (0.4604) & (0.3789) & (0.4126) & (0.3375) & (0.1759) & (0.1877) & (0.2174) & (0.2055) & (0.1937) & (0.1738) \\
Lagged\_DV & \textbf{0.7547}$^{***}$ & \textbf{0.6569}$^{***}$ & \textbf{0.7140}$^{***}$ & \textbf{0.6180}$^{***}$ & \textbf{0.7310}$^{***}$ & \textbf{0.6318}$^{***}$ & \textbf{0.7540}$^{***}$ & \textbf{0.6361}$^{***}$ & \textbf{0.7360}$^{***}$ & \textbf{0.6135}$^{***}$ & \textbf{0.7676}$^{***}$ & \textbf{0.6518}$^{***}$ \\
 & (0.0150) & (0.0186) & (0.0167) & (0.0199) & (0.0168) & (0.0208) & (0.0156) & (0.0185) & (0.0173) & (0.0205) & (0.0162) & (0.0197) \\
StockPrice &  &  & -0.0011 & -0.0033 & 0.0002 & 0.0019 &  &  & -0.0001 & \textbf{0.0101}$^{***}$ & -0.0016 & \textbf{0.0047}$^{**}$ \\
 &  &  & (0.0014) & (0.0024) & (0.0013) & (0.0024) &  &  & (0.0013) & (0.0024) & (0.0012) & (0.0022) \\
Turnover &  &  & \textbf{-24.1012}$^{***}$ & \textbf{-23.8647}$^{***}$ & \textbf{-20.7840}$^{***}$ & \textbf{-20.3411}$^{***}$ &  &  & \textbf{-7.9396}$^{***}$ & \textbf{-5.5050}$^{***}$ & \textbf{-7.1303}$^{***}$ & \textbf{-4.7073}$^{**}$ \\
 &  &  & (1.7508) & (2.2119) & (1.7290) & (2.1791) &  &  & (1.5411) & (2.0904) & (1.4699) & (1.9633) \\
DailyVola &  &  & \textbf{0.1229}$^{***}$ & \textbf{0.1393}$^{***}$ & \textbf{0.1075}$^{***}$ & \textbf{0.1257}$^{***}$ &  &  & \textbf{0.0404}$^{***}$ & \textbf{0.0520}$^{***}$ & \textbf{0.0396}$^{***}$ & \textbf{0.0441}$^{***}$ \\
 &  &  & (0.0072) & (0.0081) & (0.0082) & (0.0099) &  &  & (0.0062) & (0.0068) & (0.0069) & (0.0081) \\
AbsSurpDec &  &  & \textbf{-0.0919}$^{***}$ & -0.0241 & \textbf{-0.0756}$^{***}$ & -0.0093 &  &  & -0.0050 & 0.0119 & -0.0174 & -0.0012 \\
 &  &  & (0.0265) & (0.0300) & (0.0253) & (0.0286) &  &  & (0.0278) & (0.0304) & (0.0261) & (0.0285) \\
\midrule
Extended Controls &  &  & Yes & Yes & Yes & Yes &  &  & Yes & Yes & Yes & Yes \\
Industry FE & Yes &  & Yes &  & Yes &  & Yes &  & Yes &  & Yes &  \\
Firm FE &  & Yes &  & Yes &  & Yes &  & Yes &  & Yes &  & Yes \\
Year FE & Yes & Yes & Yes & Yes &  &  & Yes & Yes & Yes & Yes &  &  \\
Year-Quarter FE &  &  &  &  & Yes & Yes &  &  &  &  & Yes & Yes \\
\midrule
N & 42,625 & 42,625 & 42,625 & 42,625 & 42,625 & 42,625 & 43,049 & 43,049 & 43,049 & 43,049 & 43,049 & 43,049 \\
$R^2$ & 0.703 & 0.517 & 0.721 & 0.551 & 0.730 & 0.548 & 0.700 & 0.499 & 0.702 & 0.504 & 0.723 & 0.523 \\
Adj.~$R^2$ & 0.703 & 0.501 & 0.721 & 0.535 & 0.730 & 0.532 & 0.700 & 0.482 & 0.702 & 0.487 & 0.723 & 0.506 \\
\bottomrule
\end{tabular}
\begin{minipage}{\linewidth}
\vspace{2pt}\scriptsize
\textit{Notes:} $^{*}p<0.10$, $^{**}p<0.05$, $^{***}p<0.01$ (one-tailed for IVs, $\beta > 0$; two-tailed for controls).
 Significant coefficients in \textbf{bold}.
 Standard errors (in parentheses) clustered at firm level.
 Main sample (excludes financial and utility firms). DWZ Full decomposition: Clarity = -CEO FE; UncRes = call-level residual.
 $R^2$ includes absorbed fixed effects (not within-$R^2$).
 Spread$_{25D}$ values rescaled by $10^4$ (basis points) for readability; multiply coefficients by $10^{-4}$ to recover raw fraction units.
\end{minipage}
\end{table}
\clearpage

% --- tab:h18_ceo2_decomp
\begin{table}[htbp]
\centering
\caption{CEO Uncertainty and SEC Comment-Letter Receipt}
\label{tab:h18_ceo2_decomp}
\scriptsize
\begin{tabular}{lcccccc}
\toprule
 & (1) & (2) & (3) & (4) & (5) & (6) \\
 & \multicolumn{6}{c}{CCCL} \\
\cmidrule(lr){2-7}
\midrule
ClarityCEO & 0.0059 & 0.0039 & 0.0062 & 0.0040 & 0.0058 & 0.0040 \\
 & (0.0015) & (0.0030) & (0.0015) & (0.0030) & (0.0015) & (0.0030) \\
UncResCEO & -0.0003 & -0.0003 & -0.0003 & -0.0003 & -0.0004 & -0.0004 \\
 & (0.0009) & (0.0009) & (0.0009) & (0.0009) & (0.0009) & (0.0009) \\
UncPreCEO & 0.0008 & \textbf{0.0014}$^{*}$ & \textbf{0.0016}$^{**}$ & \textbf{0.0014}$^{*}$ & 0.0007 & \textbf{0.0013}$^{*}$ \\
 & (0.0008) & (0.0009) & (0.0007) & (0.0009) & (0.0008) & (0.0009) \\
\midrule
lnAssets & \textbf{-0.0006}$^{**}$ & \textbf{0.0018}$^{*}$ & 0.0001 & 0.0017 & \textbf{-0.0006}$^{**}$ & \textbf{0.0019}$^{*}$ \\
 & (0.0002) & (0.0011) & (0.0001) & (0.0010) & (0.0003) & (0.0011) \\
TobinsQ & -0.0003 & 0.0002 & -0.0001 & 0.0002 & -0.0002 & 0.0004 \\
 & (0.0003) & (0.0004) & (0.0003) & (0.0004) & (0.0003) & (0.0004) \\
ROA & 0.0006 & 0.0012 & -0.0014 & -0.0006 & -0.0006 & -0.0004 \\
 & (0.0034) & (0.0052) & (0.0052) & (0.0062) & (0.0051) & (0.0062) \\
Leverage & -0.0008 & -0.0019 & -0.0009 & -0.0011 & -0.0009 & -0.0012 \\
 & (0.0015) & (0.0034) & (0.0015) & (0.0035) & (0.0015) & (0.0035) \\
Capex & -0.0029 & 0.0106 & 0.0004 & 0.0111 & -0.0033 & 0.0107 \\
 & (0.0063) & (0.0104) & (0.0064) & (0.0104) & (0.0066) & (0.0103) \\
CashRatio & -0.0014 & \textbf{-0.0075}$^{*}$ & 0.0009 & \textbf{-0.0080}$^{*}$ & -0.0014 & \textbf{-0.0082}$^{*}$ \\
 & (0.0027) & (0.0044) & (0.0026) & (0.0044) & (0.0028) & (0.0044) \\
DivDummy & -0.0005 & 0.0003 & -0.0005 & 0.0002 & -0.0005 & 0.0002 \\
 & (0.0008) & (0.0015) & (0.0008) & (0.0015) & (0.0008) & (0.0015) \\
sCFO & -0.0007 & -0.0020 & -0.0004 & -0.0019 & -0.0007 & -0.0019 \\
 & (0.0005) & (0.0024) & (0.0004) & (0.0024) & (0.0005) & (0.0024) \\
Lagged\_DV & \textbf{-0.0079}$^{***}$ & \textbf{-0.0409}$^{***}$ & \textbf{-0.0077}$^{***}$ & \textbf{-0.0410}$^{***}$ & \textbf{-0.0074}$^{***}$ & \textbf{-0.0405}$^{***}$ \\
 & (0.0008) & (0.0027) & (0.0007) & (0.0027) & (0.0008) & (0.0028) \\
SalesGrowth &  &  & -0.0002 & -0.0013 & -0.0001 & -0.0011 \\
 &  &  & (0.0008) & (0.0011) & (0.0008) & (0.0011) \\
RDSales &  &  & -0.0003 & -0.0002 & -0.0003 & -0.0001 \\
 &  &  & (0.0002) & (0.0002) & (0.0002) & (0.0002) \\
CashFlowAt &  &  & 0.0025 & 0.0043 & 0.0009 & 0.0037 \\
 &  &  & (0.0050) & (0.0064) & (0.0050) & (0.0064) \\
DailyVola &  &  & 0.0000 & \textbf{-0.0000}$^{*}$ & 0.0000 & -0.0000 \\
 &  &  & (0.0000) & (0.0000) & (0.0000) & (0.0000) \\
\midrule
Extended Controls &  &  & Yes & Yes & Yes & Yes \\
Industry FE & Yes &  & Yes &  & Yes &  \\
Firm FE &  & Yes &  & Yes &  & Yes \\
Year FE & Yes & Yes & Yes & Yes &  &  \\
Year-Quarter FE &  &  &  &  & Yes & Yes \\
\midrule
N & 44,113 & 44,113 & 44,096 & 44,096 & 44,096 & 44,096 \\
$R^2$ & 0.001 & 0.002 & 0.000 & 0.002 & 0.001 & 0.002 \\
Adj.~$R^2$ & -0.000 & -0.032 & -0.001 & -0.032 & -0.001 & -0.033 \\
\bottomrule
\end{tabular}
\begin{minipage}{\linewidth}
\vspace{2pt}\scriptsize
\textit{Notes:} $^{*}p<0.10$, $^{**}p<0.05$, $^{***}p<0.01$ (one-tailed for IVs, $\beta > 0$; two-tailed for controls).
 Significant coefficients in \textbf{bold}.
 Standard errors (in parentheses) clustered at firm level.
 Main sample (excludes financial and utility firms). DWZ Full decomposition: Clarity = -CEO FE; UncRes = call-level residual.
 $R^2$ includes absorbed fixed effects (not within-$R^2$).
\end{minipage}
\end{table}
\clearpage


\end{document}
